\section{The algebraic induction: KT Theorem 5.5 and its three cases}
\label{sec:molecular-algebraic-induction}

This chapter proves the geometric heart of the molecular conjecture,
Katoh--Tanigawa's KT Theorem~5.5 \cite[\S5--\S6]{katohTanigawa2011}, which
turns the combinatorial reduction of \cref{sec:molecular-induction} into a
statement about rigidity-matrix rank. Where that chapter reduced every
minimal $0$-dof-graph to the two-vertex double edge by combinatorial
induction (KT Theorem~4.9, \cref{thm:minimal-kdof-reduction}), here we
\emph{realize} the reduction: every minimal $k$-dof-graph $G$ with
$|V| \ge 2$ carries a panel-hinge realization $(G,p)$ whose rigidity matrix
attains the rank $D(|V|-1) - k$ (\cref{thm:theorem-55}, at every dimension
$n \ge 3$, i.e.\ $D = \binom{n+1}{2} \ge 6$; the $d = 3$ instance is
\cref{thm:theorem-55-d3-instance}). Two consequences follow: the generic-rank
reconciliation of Tay--Whiteley's body-bar count with the rigidity-matrix
rank (KT Proposition~1.1, \cref{prop:rigidity-matrix-prop11}), and its lift
from minimal $0$-dof graphs to arbitrary deficiency (KT Theorem~5.6,
\cref{thm:theorem-55-6}, with $d = 3$ instance \cref{thm:theorem-55-6-d3}).
Together these give the molecular conjecture itself (KT Conjecture~1.2,
\cref{thm:molecular-conjecture}).

Throughout, $D = \binom{n+1}{2}$ is the body-bar dimension, $\tilde G =
(D-1)\cdot G$ is the multiplied graph of \cref{sec:body-hinge}, $R(G,p)$ is
the panel-hinge rigidity matrix of \cref{def:rigidity-matrix}, and
$\operatorname{def}(\tilde G)$ is the $D$-deficiency of
\cref{def:D-deficiency}. The rank is always read $V(G)$-relative: the
rigidity rows of $R(G,p)$ come only from the edges of $G$, so the count
depends on $G$, not on any ambient bodies outside $V(G)$. This is what lets
a realization compose through an induction that removes one vertex at a
time --- an inductive subgraph does not span the ambient body type, so its
\emph{absolute} null space always carries the free motions of the omitted
bodies (\cref{def:rank-hypothesis}, \cref{rem:rank-hypothesis-relative}).

\medskip
\noindent\textbf{The argument.} \Cref{sec:molecular-algebraic-induction-panel}
builds the \emph{panel layer}: KT's hinge-coplanar frameworks, presented in
panel-normal form so that transversality and coplanarity live in the
extensor algebra rather than in affine-subspace intersections.
\Cref{sec:molecular-algebraic-induction-prop11} proves the analytic half of
KT Proposition~1.1, identifying the graphs a generic panel realization can
make rigid. KT Theorem~5.5 is then proved by induction on $|V|$, driven by
KT Theorem~4.9 (\cref{thm:minimal-kdof-reduction}): every minimal
$k$-dof-graph on at least three vertices has a cut edge, a proper rigid
subgraph, or a splittable vertex, and each of the three moves is realized at
full rank.
\begin{itemize}
  \item \emph{Case~I} (\cref{lem:case-I}, \cref{lem:case-I-realization}), the
        contraction of a proper rigid subgraph, is KT's block-triangular
        rank-addition (KT eq.~(6.3)): the rigid-block rows and the surviving
        rows are made jointly independent by an exterior-column projection,
        reaching full rank with no shortfall.
  \item \emph{Case~II} (\cref{lem:case-II}, \cref{lem:case-II-realization})
        handles the cut edge and the positive-deficiency splitting-off, the
        panel-hinge analogue of Whiteley's bar-joint $1$-extension.
  \item \emph{Case~III} (\cref{lem:case-III}), the zero-deficiency
        splitting-off with no proper rigid subgraph, is the crux. KT's
        degenerate placement (KT eq.~(6.12)) falls one row short of full
        rank, and the missing row is supplied by a redundant-edge argument
        (KT Lemma~6.10, developed in
        \cref{sec:molecular-algebraic-induction-caseIII}).
\end{itemize}

\noindent The shared analytic engine is the genericity argument (KT
Claim~6.4, \cref{lem:genericity-device}), which selects the coordinates
attaining the generic rank; the cycle-realization lemma (KT Lemma~5.4,
\cref{lem:cycle-realization}; Crapo--Whiteley~\cite{crapoWhiteley1982},
Whiteley~\cite{whiteley1999}) enters as an input to Case~III. The induction
bottoms out on the two-vertex double edge (\cref{lem:theorem-55-base}). At
$d = 3$ the three cases assemble into \cref{thm:theorem-55-d3-instance}; the
general-dimension assembly (KT Lemma~6.13) gives \cref{thm:theorem-55}.

\subsection{The panel layer}
\label{sec:molecular-algebraic-induction-panel}

Katoh--Tanigawa's \emph{panel-hinge} framework is a
\emph{hinge-coplanar} body-hinge framework: at each body all incident
hinges lie in a common hyperplane \cite[p.\,647]{katohTanigawa2011}. The
molecular conjecture's content is that this coplanarity can be met
without losing rigidity; the Phase-18 body-hinge framework
(\cref{def:hinge-constraint}) carries free hinges with no coplanarity,
so the realization-existence nodes below need a \emph{panel layer} on
top. We take the panel-data form: a panel realization assigns each body
$v$ a hyperplane normal $n_v \in \R^{k+2}$, and the hinge at an edge
$e = uv$ is the codimension-2 intersection $\mathrm{panel}(u) \cap
\mathrm{panel}(v)$. Its supporting $k$-extensor --- the screw-space
element of \cref{def:hinge-constraint} the rigidity matrix constrains ---
is the Grassmann--Cayley meet of the two panels, equivalently the
complement iso of the join of the two normals. Coplanarity (both hinges
at $v$ lie in $\mathrm{panel}(v)$ by construction) and transversality
both live in the extensor algebra, so no affine-subspace-intersection
plumbing is needed.

\begin{definition}[Panel support extensor]
  \label{def:panel-support-extensor}
  \lean{CombinatorialRigidity.Molecular.panelSupportExtensor,
        CombinatorialRigidity.Molecular.normalsJoin,
        CombinatorialRigidity.Molecular.panelSupportExtensor_ne_zero_iff}
  \leanok
  \uses{def:join, def:meet-complement-iso, def:hinge-constraint}
  For two panel normals $n_1, n_2 \in \R^{k+2}$, the \emph{panel support
  extensor} is the supporting $k$-extensor $C(p(e)) \in \Lambda^k
  \R^{k+2}$ of the codimension-2 intersection of the two panels, given as
  the complement iso (\cref{def:meet-complement-iso}) of the grade-2 join
  $n_1 \wedge n_2 \in \Lambda^2 \R^{k+2}$ (\cref{def:join}). Since
  $k+2-2 = k$ is the screw-space grade, the result lands in the screw
  space $\Lambda^k \R^{k+2}$ of \cref{def:hinge-constraint} and feeds the
  hinge constraint verbatim. It is nonzero exactly when the two normals
  are linearly independent --- the algebraic form of the two panels
  meeting transversally in a genuine hinge --- the only general-position
  condition the panel layer needs.
\end{definition}

\begin{definition}[Panel-hinge framework and its body-hinge interpretation]
  \label{def:panel-hinge-framework}
  \lean{CombinatorialRigidity.Molecular.PanelHingeFramework,
        CombinatorialRigidity.Molecular.PanelHingeFramework.toBodyHinge,
        CombinatorialRigidity.Molecular.BodyHingeFramework.IsHingeCoplanar,
        CombinatorialRigidity.Molecular.BodyHingeFramework.isHingeCoplanar_toBodyHinge,
        CombinatorialRigidity.Molecular.PanelHingeFramework.IsGeneralPosition,
        CombinatorialRigidity.Molecular.PanelHingeFramework.supportExtensor_ne_zero_of_isGeneralPosition,
        CombinatorialRigidity.Molecular.PanelHingeFramework.isGeneralPosition_withMomentNormals,
        CombinatorialRigidity.Molecular.PanelHingeFramework.isGeneralPosition_ofParam}
  \leanok
  \uses{def:panel-support-extensor, def:hinge-constraint}
  A \emph{panel-hinge framework} is a multigraph $G = (V,E)$ together with a
  per-body panel normal $n_v \in \R^{k+2}$ (the pole of body $v$'s
  hyperplane $\mathrm{panel}(v)$). The hinge at an edge $e = uv$ is the
  codimension-2 intersection $\mathrm{panel}(u) \cap \mathrm{panel}(v)$,
  with supporting $k$-extensor the meet $C(p(e)) =
  \mathrm{panelSupportExtensor}(n_u, n_v)$
  (\cref{def:panel-support-extensor}). Its \emph{body-hinge interpretation}
  is the body-hinge framework (\cref{def:hinge-constraint}) on $G$ carrying
  exactly these supporting extensors, so the entire rigidity-matrix rank
  theory of \cref{sec:molecular-rigidity-matrix} applies verbatim. Every
  hinge incident to $v$ has $n_v$ as a join factor, hence lies in the single
  panel $\mathrm{panel}(v)$: the framework is \emph{hinge-coplanar} by
  construction. A body-hinge framework is hinge-coplanar
  ($\mathrm{IsHingeCoplanar}$) when it arises this way; this is the
  molecular-model constraint Katoh--Tanigawa's panel-hinge frameworks carry
  beyond the free-hinge body-hinge model, and the conjecture's content is
  that it can be met without dropping rigidity. The normals are in
  \emph{general position} ($\mathrm{IsGeneralPosition}$) when any two normals
  at distinct bodies are linearly independent; then every hinge joining two
  distinct bodies has a nonzero supporting extensor
  ($\mathrm{supportExtensor\_ne\_zero\_of\_isGeneralPosition}$) --- the
  transversality the rigid-block forest hinges feed into the block-triangular
  gluing of \cref{lem:case-I}. An explicit, dimension-free general-position
  assignment exists for any number of bodies: setting each body's normal to a
  point $(1, t, t^2, \dots, t^{k+1})$ on the moment curve at a distinct
  parameter $t$ makes any two normals at distinct bodies independent (the
  $2 \times 2$ Vandermonde minor has determinant $t' - t \neq 0$), so an
  injective real assignment of parameters yields
  $\mathrm{IsGeneralPosition}$
  ($\mathrm{isGeneralPosition\_withMomentNormals}$, and from a bare graph and
  endpoint selector directly via $\mathrm{isGeneralPosition\_ofParam}$ --- the
  realization-side entry point for the Case-I assembly)
  --- where standard-basis normals would cover only $|V| \le k+2$.
\end{definition}

\subsection{The generic-rank reconciliation (Proposition 1.1, analytic half)}
\label{sec:molecular-algebraic-induction-prop11}

The matroidal half of Tay--Whiteley's Proposition~1.1 --- the bridge
$\operatorname{def}(\tilde G) = \operatorname{corank} M(\tilde G)$ ---
landed green in Phase~19 (\cref{thm:def-eq-corank}). The analytic half
relocated forward to this phase: it wires the rigidity-matrix rank
$R(G,p)$ of \cref{sec:molecular-rigidity-matrix} to the matroid corank,
and depends on the generic-rank argument (Claim~6.4) that drives the
whole induction.

\begin{proposition}[Generic rank equals $D(|V|-1) - \operatorname{def}$; KT Proposition 1.1, analytic half]
  \label{prop:rigidity-matrix-prop11}
  \lean{CombinatorialRigidity.Molecular.rigidityMatrix_prop11}
  \leanok
  \uses{def:rigidity-matrix, def:D-deficiency, thm:def-eq-corank, def:dof-generic, lem:genericity-device}
  For a generic panel-hinge realization $(G,p)$,
  \[
    \operatorname{rank} R(G,p) = D(|V|-1) - \operatorname{def}(\tilde G).
  \]
  Equivalently $G$ has a rigid panel-hinge realization iff
  $\tilde G = (D-1)\cdot G$ contains $D$ edge-disjoint spanning trees,
  reconciling the honest rank form of
  \cref{sec:molecular-rigidity-matrix} with the body-hinge existence
  form (\cref{thm:body-hinge-tay}) of Phase~16. This is
  Jackson--Jord\'an's Theorem~6.1, geometric side
  \cite{jacksonJordan2009}.
\end{proposition}
\begin{proof}
  \leanok
  \uses{thm:theorem-55, lem:trivial-motions-rank-bound,
        lem:motions-mono-of-graph-le}
  The deficiency is the corank of $M(\tilde G)$
  (\cref{thm:def-eq-corank}), so the claim is that a generic realization
  attains the maximum rank $D(|V|-1) - \operatorname{corank}$. The target
  equality is the realization hypothesis \cref{def:rank-hypothesis} at
  $k' = \operatorname{def}$; it is pinned by two inequalities. The
  genericity-free upper bound $\operatorname{rank} R(G,p) \le D(|V|-1) -
  \operatorname{def}$ is the codimension form
  \cref{lem:trivial-motions-rank-bound} together with the deficiency count: a
  partition $P$ attaining $\operatorname{def}(\tilde G)$ contributes
  $\operatorname{def}$ independent motions beyond the $D$ trivial ones --- this
  holds for \emph{every} realization, no genericity, and is discharged in-proof
  from the Phase-19 partition machinery. The reverse (generic max-rank) bound
  $\operatorname{rank} R(G,p) \ge D(|V|-1) - \operatorname{def}$ --- equivalently
  $\dim Z(G,p) \le D + \operatorname{def}$, carried as the hypothesis
  $\mathtt{hgen}$ --- holds at a generic realization: delete $G$ down to a
  minimal $k$-dof spanning subgraph, realize it at full generic rank by
  Theorem~5.5 (\cref{thm:theorem-55}, the Claim~6.4 generic-rank argument
  selecting the attaining point), then re-add the deleted edges, which only
  grows the rank (\cref{lem:motions-mono-of-graph-le}). The proposition's only
  further inputs are the dimension fixing $n = k+1$ and the genuine-hinge
  condition $C(e) \neq 0$ the partition cut needs. The construction supplying
  $\mathtt{hgen}$ at every deficiency is Theorem~5.6 at $d = 3$
  (\cref{thm:theorem-55-6-d3}). Applied to a spanning $G$ ($V(G) = V_\alpha$),
  the $V(G)$-relative motive of \cref{def:rank-hypothesis} coincides with the
  absolute $\dim Z = D + \operatorname{def}$.
\end{proof}

\subsection{The induction skeleton and base case}
\label{sec:molecular-algebraic-induction-base}

\begin{definition}[Generic-rank hypothesis (6.1), $V(G)$-relative --- rank-form substrate]
  \label{def:rank-hypothesis}
  \lean{CombinatorialRigidity.Molecular.BodyHingeFramework.IsInfinitesimallyRigidOn,
        CombinatorialRigidity.Molecular.PanelHingeFramework.HasFullRankRealization,
        CombinatorialRigidity.Molecular.PanelHingeFramework.HasGenericFullRankRealization,
        CombinatorialRigidity.Molecular.BodyHingeFramework.RankHypothesis}
  \leanok
  \uses{def:rigidity-matrix, def:k-dof}
  For a minimal $k$-dof-graph $G$ with $|V| \ge 2$, the
  \emph{realization hypothesis} $(6.1)$ asserts the existence of a
  panel-hinge realization $(G,p)$ with
  $\operatorname{rank} R(G,p) = D(|V|-1) - k$.

  \emph{The rank must be read $V(G)$-relative.} The rigidity rows of
  $R(G,p)$ come only from the edges of $G$, so $\operatorname{rank} R(G,p) =
  \dim \operatorname{span}(\text{rigidity rows})$ depends only on $V(G)$ ---
  it is blind to any ambient bodies not in $V(G)$. The target form the Lean
  carries is therefore the \emph{rank} equality
  \[
    \dim \operatorname{span}\bigl(\text{rigidity rows of } (G,p)\bigr)
      \;=\; D\,(|V(G)| - 1) - \operatorname{def}(\tilde{G}),
  \]
  equivalently ``every infinitesimal motion is constant on $V(G)$'' when
  $\operatorname{def} = 0$. This is \emph{not} the same as the null-space equality
  $\dim Z(G,p) = D + k$ over the whole screw-assignment space
  $\R^{D|V_\alpha|}$ of the ambient body type; the rank form above is the
  $\alpha$-independent motive the realization layer is stated against.

  \textbf{Lean forms.}  Two predicates pin this definition.

  $\mathtt{HasFullRankRealization}\,k\,G$ is the bare existential
  $\exists\,Q,\; Q.\mathrm{graph} = G \;\wedge\;
  Q.\mathrm{toBodyHinge}.\mathrm{IsInfinitesimallyRigidOn}\,V(G)$,
  used as the induction motive of Theorem~5.5.
  It is equivalent to the rank equality (B1,
  $\mathtt{isInfinitesimallyRigidOn\_vertexSet\_iff\_finrank\_span\_rigidityRows}$)
  when $V(G)$ is nonempty.  It remains \emph{weaker} than KT's (6.1) ---
  a degenerate welded framework satisfies it on every connected $0$-dof
  graph --- and is slated for deletion once the spine restates onto the
  honest genuine-hinge motive (\cref{def:genuine-hinge-realization}).

  $\mathtt{HasGenericFullRankRealization}\,k\,n\,G$ is the strengthened motive
  that carries the rank equality as an explicit conjunct (at deficiency $n$,
  alongside general position and algebraic independence of the normals);
  it is the induction substrate for the generic-motive reduction
  $\mathtt{theorem\_55\_generic}$.  The rank conjunct is the
  $\mathbb{Z}$-cast form
  $(\dim \operatorname{span}(\text{rigidity rows}) : \mathbb{Z}) =
    D\,(|V(G)| - 1) - \operatorname{def}(\tilde{G})$,
  matching the $\mathbb{Z}$ arithmetic of the genuine-hinge motive
  (\cref{def:genuine-hinge-realization}).
\end{definition}

\begin{definition}[Generic-rank hypothesis (6.1) --- genuine-hinge form; KT eq.\ (6.1)]
  \label{def:genuine-hinge-realization}
  \lean{CombinatorialRigidity.Molecular.HasPanelRealization,
    CombinatorialRigidity.Molecular.ExtensorInPanel,
    CombinatorialRigidity.Molecular.hasPanelRealization_of_generic}
  \leanok
  \uses{def:rank-hypothesis, def:panel-hinge-framework, def:D-deficiency}
  A graph $G$ \emph{has a genuine-hinge panel realization} at dimension $k$ and
  defect parameter $n$ when there exist a body-hinge framework $F$ on $G$ and a
  panel-normal assignment $\mathtt{normal} : V \to \mathbb{R}^{k+2}$ such that:
  every vertex has a nonzero normal; every link's supporting $k$-extensor is
  nonzero and lies in both endpoint panels (each endpoint's hyperplane
  $\mathtt{normal}(v)^\perp$ contains $k$ spanning points of the extensor,
  $\mathtt{ExtensorInPanel}$); and the rigidity-row span attains the target rank
  $D(|V(G)|-1) - \operatorname{def}(\tilde G)$ (KT eq.~(6.1), $\mathbb Z$-cast form).

  Every generic realization (\cref{def:panel-hinge-framework}) with
  $|V| \ge 2$ yields a genuine-hinge realization
  ($\mathtt{hasPanelRealization\_of\_generic}$): panel normals in general position
  give nonzero extensors (general position forces transversality on each link);
  the meet-decomposition lemma furnishes the $\mathtt{ExtensorInPanel}$ witnesses;
  and an infinitesimally rigid framework attains the rank count exactly.
\end{definition}

\begin{theorem}[Realization at the target rank; KT Theorem 5.5]
  \label{thm:theorem-55}
  \lean{CombinatorialRigidity.Molecular.PanelHingeFramework.theorem_55_gen,
        CombinatorialRigidity.Molecular.PanelHingeFramework.theorem_55_minimalKDof_gen}
  \leanok
  \uses{def:rank-hypothesis, thm:minimal-kdof-reduction, lem:case-I,
        lem:case-II, lem:case-I-realization, lem:case-I-dispatch,
        lem:case-II-realization, lem:case-III, lem:theorem-55-base,
        lem:case-III-nested-rank-lower}
  Every minimal $0$-dof-graph $G$ with $|V| \ge 2$ has a panel-hinge
  realization $(G,p)$ with $\operatorname{rank} R(G,p) = D(|V(G)|-1)$ ---
  an infinitesimally rigid (on $V(G)$) panel-hinge framework
  (\cref{def:rank-hypothesis}, at $k' = 0$). Formalized at general $d$ for
  $D = \binom{n+1}{2} \ge 6$, i.e. $n \ge 3$ (the Phase-20 chain extractors
  fix this floor); a fresh-edge supply rides as ambient data. (The molecular
  conjecture needs only the $0$-dof instance stated here. The proof is not
  self-contained in $k$: Katoh--Tanigawa's induction runs over minimal
  $k$-dof-graphs for \emph{all} $k$, and the Case~III argument applies the
  inductive statement to an auxiliary minimal $k'$-dof graph with
  $0 \le k' \le D-2$ (Claim~6.11, eq.~(6.22)), discharged from the all-$k$ IH
  (\cref{lem:case-III-nested-rank-lower}).)
\end{theorem}
\begin{proof}
  \leanok
  \uses{thm:minimal-kdof-reduction, lem:theorem-55-base, lem:case-I,
        lem:case-II, lem:case-I-realization, lem:case-I-dispatch,
        lem:case-II-realization, lem:case-III, lem:case-III-nested-rank-lower,
        lem:chain-data-extract}
  Induction on $|V(G)|$, driven by the combinatorial reduction
  Theorem~4.9 (\cref{thm:minimal-kdof-reduction}): the all-$k$ well-founded
  induction principle (\texttt{minimal\_kdof\_reduction\_all\_k}) is run
  against the simplicity-conditioned motive
  ``$(G.\mathrm{Simple} \Rightarrow \mathrm{GP}) \wedge
    \mathrm{HasPanelRealization}$'' (\cref{def:rank-hypothesis}). The base
  case $|V(G)| = 2$ is \cref{lem:theorem-55-base}. The inductive step splits
  on: a cut edge (Case~0, \cref{lem:case-II-realization}); a proper rigid
  subgraph (Case~I, \cref{lem:case-I-dispatch}); $k > 0$ with no rigid
  subgraph (Case~II, \cref{lem:case-II-realization}); and $k = 0$ with no
  rigid subgraph (Case~III, \cref{lem:case-III}). The nested all-$k$ IH
  carries through all branches; Case~III's nested rank bound is
  \cref{lem:case-III-nested-rank-lower}. The grade-general spine
  \texttt{theorem\_55\_minimalKDof\_gen} fills the four reduction-case
  producers from their grade-general forms (base, cut, Case-I contract,
  and the forgetful map), with the Case-III arm run at general $n$ by the
  chain-dispatch router inside \texttt{case\_III\_realization\_all\_k} and
  the general chain extractor (\cref{lem:chain-data-extract}); the $d = 3$
  specialization is \cref{thm:theorem-55-d3-instance}.
\end{proof}

\begin{theorem}[Theorem 5.5 at $d = 3$: the zero-carry instance; KT Theorem 5.5, \S6.4.1]
  \label{thm:theorem-55-d3-instance}
  \lean{CombinatorialRigidity.Molecular.PanelHingeFramework.theorem_55_all_k,
        CombinatorialRigidity.Molecular.theorem_55_base_producer,
        CombinatorialRigidity.Molecular.PanelHingeFramework.theorem_55_d3,
        Graph.not_simple_of_isMinimalKDof_of_ncard_two,
        CombinatorialRigidity.Molecular.PanelHingeFramework.rankHypothesis_deficiency_of_theorem_55_d3}
  \leanok
  \uses{def:rank-hypothesis, def:genuine-hinge-realization,
        thm:minimal-kdof-reduction, lem:theorem-55-base-producer,
        lem:case-I-realization, lem:case-I-dispatch,
        lem:case-II-realization, lem:case-III, lem:case-III-nested-rank-lower}
  The $d = 3$ ($D = 6$) instance of \cref{thm:theorem-55}, run against the
  simplicity-conditioned motive: every minimal $0$-dof-graph $G$ with
  $|V| \ge 2$ satisfies the pair --- if $G$ is simple, a general-position,
  link-recording realization at an algebraically-independent seed of rank
  $D(|V|-1)$; and in all cases a genuine-hinge panel realization at rank
  $D(|V|-1)$ (\cref{def:genuine-hinge-realization}). The all-$k$ recursion
  (\texttt{minimal\_kdof\_reduction\_all\_k}) is run against the conditioned
  motive; all three adjudicated carries ($\mathtt{h622}$, $\mathtt{hsplit}$,
  $\mathtt{hcontract}$) are now discharged (Phase~22k~L7--L9): $\mathtt{h622}$
  via \cref{lem:case-III-nested-rank-lower}; $\mathtt{hsplit}$ by
  \cref{lem:simple-minimal-noRigid} (G0) $+$ Case-III
  (\cref{lem:case-III}); $\mathtt{hcontract}$ by \cref{lem:case-I-dispatch}.
  A fresh-edge supply rides as ambient data.

  \emph{The first push toward Theorem 5.6.} For a \emph{spanning} simple
  minimal $0$-dof-graph ($V(G)$ the whole body type), the general-position
  conjunct yields a panel-hinge framework realizing the rank hypothesis at
  $\operatorname{def}(\tilde G) = 0$ --- the $\operatorname{def} = 0$ generic-rank
  feed of the Proposition-1.1 reconciliation (\cref{prop:rigidity-matrix-prop11}).
  The $\operatorname{def} > 0$ form is the next theorem
  (\cref{thm:theorem-55-6-d3}).
\end{theorem}
\begin{proof}
  \leanok
  \uses{thm:minimal-kdof-reduction, lem:theorem-55-base-producer,
        lem:case-I-realization, lem:case-I-dispatch,
        lem:case-II-realization, lem:case-III, lem:case-III-nested-rank-lower,
        prop:rigidity-matrix-prop11}
  Instantiate \texttt{theorem\_55\_all\_k} at $D = 6$. The two-vertex base is
  \cref{lem:theorem-55-base-producer}. The no-rigid-subgraph branches
  ($k = 0$: \cref{lem:case-III}; $k > 0$: \cref{lem:case-II-realization}) each
  start with G0 to get $G.\mathrm{Simple}$, then apply the GP producer and
  M4 for the bare conjunct. The contraction branch dispatches on $k$:
  at $k = 0$, use \cref{lem:case-I-dispatch}; at $k > 0$, the 6.5 arm is
  vacuous (the carrier construction forces $k = 0$), and the remaining arms
  use \cref{lem:case-I-realization}. For the spanning stratum corollary,
  re-aim the endpoint selector at a fixed distinct pair on non-edges and apply
  the Proposition-1.1 rank count (\cref{prop:rigidity-matrix-prop11}).
\end{proof}

\begin{theorem}[Theorem 5.6 at $d = 3$: every multigraph attains the deficiency rank; KT Theorem 5.6, \S5.2]
  \label{thm:theorem-55-6-d3}
  \lean{CombinatorialRigidity.Molecular.PanelHingeFramework.rankHypothesis_of_theorem_55_d3,
        CombinatorialRigidity.Molecular.PanelHingeFramework.rankHypothesis_deficiency_of_theorem_55_d3}
  \leanok
  \uses{thm:theorem-55-d3-instance, lem:motions-mono-of-graph-le,
        lem:subgraph-minimality, def:rank-hypothesis}
  At $d = 3$ ($D = 6$), every simple spanning multigraph $G$ has a
  panel-hinge realization $(G,p)$ attaining the deficiency rank
  \[
    \operatorname{rank} R(G,p) = D(|V|-1) - \operatorname{def}(\tilde G),
  \]
  i.e.\ a framework realizing the rank hypothesis
  (\cref{def:rank-hypothesis}) at $k' = \operatorname{def}(\tilde G)$. This is
  Katoh--Tanigawa's Theorem~5.6 specialized to $d = 3$, lifting the
  minimal-$0$-dof realization of \cref{thm:theorem-55-d3-instance} to an
  arbitrary deficiency. It is the genuine $\operatorname{def} > 0$ feed of the
  Proposition-1.1 reconciliation (\cref{prop:rigidity-matrix-prop11}); the
  $\operatorname{def} = 0$ stratum is \cref{thm:theorem-55-d3-instance}.
\end{theorem}
\begin{proof}
  \leanok
  \uses{thm:theorem-55-d3-instance, lem:motions-mono-of-graph-le,
        lem:subgraph-minimality, prop:rigidity-matrix-prop11}
  Following Katoh--Tanigawa (p.~670). When $|V| = 1$ the framework is trivially
  rigid and $\operatorname{def}(\tilde G) = 0$, settled directly. Otherwise:
  delete edges of $G$ down to a minimal $k$-dof spanning subgraph $G' \le G$
  with $V(G') = V(G)$ and $\operatorname{def}(\tilde{G'}) =
  \operatorname{def}(\tilde G) =: k$ (the deficiency-preserving strip; the
  matroid engine behind \cref{lem:subgraph-minimality}). Realize $G'$ at its own
  deficiency $k$ by Theorem~5.5 (\cref{thm:theorem-55-d3-instance}, the all-$k$
  spine: $G'$ is simple as a subgraph of a simple $G$), yielding a generic
  full-rank realization $q$ on $G'$ with $\dim Z(G',q) = D + k$. Re-aim $q$ to
  the full graph $G$, keeping the supporting extensors on every $G'$-link
  unchanged and placing a genuine in-panel hinge on each re-added edge --- in
  the homogeneous model two distinct hyperplanes through the origin always meet,
  so KT's projective-invariance move is free. Re-adding edges only shrinks the
  motion space (\cref{lem:motions-mono-of-graph-le}), so
  $\dim Z(G,p) \le \dim Z(G',q) = D + k = D + \operatorname{def}(\tilde G)$. This
  is exactly the upper bound $\mathtt{hgen}$, which together with the
  genericity-free lower bound discharges the rank hypothesis via the
  Proposition-1.1 rank bridge (\cref{prop:rigidity-matrix-prop11}).
\end{proof}

\begin{lemma}[Theorem 5.5 base case ($|V(G)| = 2$)]
  \label{lem:theorem-55-base}
  \lean{CombinatorialRigidity.Molecular.BodyHingeFramework.theorem_55_base}
  \leanok
  \uses{def:rank-hypothesis, lem:rank-parallel-full}
  The two-vertex double edge has a panel-hinge realization of rank
  $D(|V(G)|-1) - k = D - k$. For the minimal $0$-dof case (the double edge,
  $k = 0$) this is the full rank $D$, attained by two non-parallel hinges
  via the parallel-hinges-full Lemma~5.3 (\cref{lem:rank-parallel-full}).
  At $k = 0$, in the $V(G)$-relative motive (\cref{def:rank-hypothesis}),
  the conclusion is exactly ``every infinitesimal motion is constant on the
  two bodies $V(G) = \{u, v\}$'' ($\mathrm{IsInfinitesimallyRigidOn}\,\{u,v\}$).
  The prior Lean required the constancy to hold over the \emph{whole} ambient
  body type (a hypothesis $\forall w,\ w = u \lor w = v$, i.e.\ ``$\alpha =
  \{u,v\}$''), the absolute-motive artefact; against the relativized motive the
  constancy is only required on $V(G)$, so that ambient hypothesis is
  dropped and the base case holds for any ambient type.
\end{lemma}
\begin{proof}
  \leanok
  \uses{lem:rank-parallel-full}
  Two bodies joined by hinges with linearly independent supporting
  extensors realize the full screw-difference constraint: the
  parallel-hinges-full Lemma~5.3 (\cref{lem:rank-parallel-full}) forces
  $S(u) = S(v)$ for every infinitesimal motion $S$. That is exactly
  constancy on $V(G) = \{u, v\}$, so the rank of the two hinge-row blocks is
  the full $D = D(2-1)$ and the $V(G)$-relative motive holds at $k = 0$. No
  condition on bodies outside $\{u, v\}$ is needed.
\end{proof}

\begin{lemma}[Two independent extensors in a common panel]
  \label{lem:extensor-pair-in-panel}
  \lean{CombinatorialRigidity.Molecular.exists_linearIndependent_extensor_pair_perp}
  \leanok
  \uses{def:extensor, lem:extensor-independence}
  For any normal $n \in \R^4$ there are two point-pairs $p, q : \{0,1\} \to \R^4$, each lying
  in the panel $n^{\perp} = \{x : x \cdot n = 0\}$, whose grade-2 extensors
  $\mathtt{extensor}\,p,\ \mathtt{extensor}\,q \in \bigwedge^2 \R^4$ are linearly independent.
  This is the coincident-panel ($\Pi(u) = \Pi(v) = n^{\perp}$) geometric brick of
  Katoh--Tanigawa's parallel-pair realization (Lemma~5.3, KT~2011 p.~670): two hinges sharing a
  panel but carrying non-proportional supporting extensors, the input the two-vertex base producer
  feeds to \cref{lem:theorem-55-base}.
\end{lemma}
\begin{proof}
  \leanok
  \uses{lem:extensor-independence}
  The panel $n^{\perp}$ has dimension at least $3$ in $\R^4$ (it is the kernel of the single
  pairing functional $x \mapsto x \cdot n$, so $\dim n^{\perp} \ge 4 - 1 = 3$, with equality
  unless $n = 0$). Pick three linearly independent vectors $a, b, c \in n^{\perp}$ and set
  $p = (a, b)$, $q = (a, c)$. The two wedges $a \wedge b$ and $a \wedge c$ are linearly
  independent: a relation $s\,(a \wedge b) + t\,(a \wedge c) = 0$, joined on the left with $c$
  (resp.\ $b$), kills the second (resp.\ first) term by the alternating property and leaves
  $s\,(c \wedge a \wedge b) = 0$ (resp.\ $t\,(b \wedge a \wedge c) = 0$); the surviving triple
  wedge is nonzero by independence of $\{a, b, c\}$, so $s = t = 0$. (This is the two-subset
  instance of the join-to-isolate technique behind \cref{lem:extensor-independence}.)
\end{proof}

\begin{lemma}[Theorem 5.5 base producer, empty arm]
  \label{lem:theorem-55-base-producer-empty}
  \lean{CombinatorialRigidity.Molecular.theorem_55_base_producer_empty,
        CombinatorialRigidity.Molecular.theorem_55_base_producer_empty_gp}
  \leanok
  \uses{def:rank-hypothesis, def:genuine-hinge-realization, lem:two-vertex-trichotomy}
  Let $G$ be a minimal-$k$-dof-graph on $1 \le |V(G)| \le 2$ bodies realizing the \emph{empty}
  case of the base trichotomy (\cref{lem:two-vertex-trichotomy} (i)): $E(G) = \emptyset$ and
  hence $\mathrm{def}(\tilde G) = D(|V(G)| - 1)$. Then $G$ has a genuine-hinge panel realization
  (\cref{def:genuine-hinge-realization}) at the target rank $D(|V(G)|-1) - \mathrm{def} = 0$, and
  --- the general-position companion --- a \emph{generic} full-rank realization
  (\cref{def:rank-hypothesis}) at that same rank.
\end{lemma}
\begin{proof}
  \leanok
  The bare construction is geometry-free: take the all-zero framework (every edge carries the
  zero supporting extensor) at a fixed nonzero panel normal. With no edges no hinge constraint
  fires, so the rigidity-row family is empty, its span is $\{0\}$, and the rank is $0$; the
  per-link conjunct of the motive is vacuous. The target rank vanishes because the empty-graph
  deficiency is exactly $D(|V(G)|-1)$ (\cref{lem:two-vertex-trichotomy}). For the generic
  companion build the single-body / empty framework at an injective
  algebraically-independent seed (a non-root of the general-position polynomial): the panel
  normals are then in general position and algebraically independent, the rigidity-row family is
  still empty (rank $0$), and link-recording is vacuous.
\end{proof}

\begin{lemma}[Theorem 5.5 base producer, single-edge arm]
  \label{lem:theorem-55-base-producer-single}
  \lean{CombinatorialRigidity.Molecular.theorem_55_base_producer_single_edge,
        CombinatorialRigidity.Molecular.theorem_55_base_producer_single_edge_gp}
  \leanok
  \uses{def:rank-hypothesis, def:genuine-hinge-realization, lem:two-vertex-trichotomy,
        lem:extensor-pair-in-panel, def:rigidity-matrix}
  Let $G$ be a two-vertex minimal-$1$-dof-graph realizing the \emph{single-edge} case of the base
  trichotomy (\cref{lem:two-vertex-trichotomy} (ii)): $V(G) = \{x, y\}$ with $x \ne y$,
  $E(G) = \{e\}$ a single edge joining $x$ and $y$, and $\mathrm{def}(\tilde G) = 1$. Then $G$ has
  a genuine-hinge panel realization (\cref{def:genuine-hinge-realization}) at the target rank
  $D(|V(G)|-1) - \mathrm{def} = D - 1 = 5$ (at $d = 3$, $D = 6$), and --- the general-position
  companion, the one base arm where this conjunct does genuine geometric work --- a \emph{generic}
  full-rank realization (\cref{def:rank-hypothesis}) at that same rank.
\end{lemma}
\begin{proof}
  \leanok
  \uses{lem:extensor-pair-in-panel, def:rigidity-matrix}
  For the bare motive, place coincident panels $\Pi(x) = \Pi(y) = n^{\perp}$ at a fixed nonzero
  normal $n$ and assign the single hinge one nonzero supporting extensor inside the common panel
  (the first of the pair of \cref{lem:extensor-pair-in-panel}; nonzero by its linear
  independence). The lone hinge's row block already \emph{is} the entire rigidity-row family
  (there is only one edge), and that block has rank exactly $D - 1$ (\cref{def:rigidity-matrix});
  so the rigidity-row span has rank $D - 1 = D(|V(G)|-1) - \mathrm{def}$, the rank equality of the
  motive. Unlike the parallel-pair arm the framework is not rigid ($\mathrm{def} = 1$): the rank
  is the single hinge-row count, not the full two-vertex base rank. The generic companion builds
  $\mathtt{ofNormals}$ at an injective algebraically-independent seed (a non-root of the
  general-position polynomial) with the single edge's two distinct endpoints: general position
  forces the lone supporting extensor nonzero, the same single hinge-row count gives rank $D - 1$,
  and the seed supplies general position $+$ algebraic independence of the panel normals.
\end{proof}

\begin{lemma}[Theorem 5.5 base producer, parallel-pair arm]
  \label{lem:theorem-55-base-producer-parallel}
  \lean{CombinatorialRigidity.Molecular.theorem_55_base_producer_parallel_pair}
  \leanok
  \uses{def:genuine-hinge-realization, lem:two-vertex-trichotomy, lem:theorem-55-base,
        lem:extensor-pair-in-panel}
  Let $G$ be a two-vertex minimal-$0$-dof-graph realizing the \emph{parallel-pair} case of the
  base trichotomy (\cref{lem:two-vertex-trichotomy} (iii)): $V(G) = \{x, y\}$ with $x \ne y$,
  $E(G) = \{e, f\}$ with $e \ne f$ two edges both joining $x$ and $y$, and $\mathrm{def}(\tilde G)
  = 0$. Then $G$ has a genuine-hinge panel realization
  (\cref{def:genuine-hinge-realization}) at the full target rank $D(|V(G)|-1) - \mathrm{def} =
  D \cdot 1 = 6$ (at $d = 3$, $D = 6$). This is the geometrically nontrivial arm of the all-$k$
  base producer; the empty and single-edge arms are bookkeeping / single-row counts.
\end{lemma}
\begin{proof}
  \leanok
  \uses{lem:theorem-55-base, lem:extensor-pair-in-panel}
  Place \emph{coincident panels} $\Pi(x) = \Pi(y) = n^{\perp}$ at a fixed nonzero normal
  $n$ and assign the two parallel hinges two linearly independent supporting extensors inside the
  common panel $n^{\perp}$ (\cref{lem:extensor-pair-in-panel}). Each link's supporting extensor is
  then nonzero and lies in both endpoint panels, discharging the per-link conjunct of the motive.
  The two independent extensors give the combined hinge-row blocks full rank $D$ on the relative
  screw $S(x) - S(y)$, so the two-vertex base case (\cref{lem:theorem-55-base}) makes the framework
  infinitesimally rigid on $\{x, y\} = V(G)$; the $V(G)$-relative rank bridge (rigid on $V(G)$ iff
  the rigidity-row span has rank $D(|V(G)|-1)$) turns that into the rank equality of the motive,
  with $\mathrm{def}(\tilde G) = 0$.
\end{proof}

\begin{lemma}[Theorem 5.5 base producer, trichotomy dispatch]
  \label{lem:theorem-55-base-producer}
  \lean{CombinatorialRigidity.Molecular.theorem_55_base_producer}
  \leanok
  \uses{def:rank-hypothesis, def:genuine-hinge-realization, lem:two-vertex-trichotomy,
        lem:theorem-55-base-producer-empty, lem:theorem-55-base-producer-single,
        lem:theorem-55-base-producer-parallel}
  Let $G$ be a minimal-$k$-dof-graph with $1 \le |V(G)| \le 2$ (the base region of the
  all-$k$ reduction, \cref{thm:minimal-kdof-reduction-all-k}).  Then the conditioned pair
  $(G\text{ simple} \Rightarrow \text{has a \emph{generic} full-rank panel realization})$
  and $(\text{has a genuine-hinge panel realization at target rank})$ both hold.  The
  simple conjunct is the strengthened generic motive
  (\cref{def:rank-hypothesis}, general position $+$ algebraically-independent panel
  normals); the unconditional conjunct is the honest bare motive
  (\cref{def:genuine-hinge-realization}).  This is the conditioned pair $P_c\,G$ that the
  all-$k$ spine recurses against.
\end{lemma}
\begin{proof}
  \leanok
  \uses{lem:theorem-55-base-producer-empty, lem:theorem-55-base-producer-single,
        lem:theorem-55-base-producer-parallel}
  Dispatch on the base trichotomy (\cref{lem:two-vertex-trichotomy}). The unconditional
  bare conjunct comes from the three bare arms. For the simple conjunct: the empty arm
  ($E(G) = \emptyset$) and the single-edge arm ($\mathrm{def} = 1 > 0$) both admit a
  \emph{generic} full-rank realization --- built from $\mathtt{ofNormals}$ at an injective
  algebraically-independent seed, which is a non-root of the general-position polynomial, so
  the panel normals are simultaneously in general position and algebraically independent; the
  rank is $0$ (empty) resp.\ $D - 1$ (single-edge, the single hinge-row count). The
  single-edge arm is the one base arm where the simple conjunct does genuine geometric work.
  In the parallel-pair arm ($k = 0$) a two-vertex minimal-$0$-dof-graph cannot be simple (the
  single-vertex cut supplies two distinct parallel edges, both forced to link the same pair,
  which simplicity collapses --- contradiction); the $G$-simple antecedent is therefore
  vacuous.
\end{proof}

\begin{lemma}[Theorem 5.5 triangle base ($|V(G)| = 3$)]
  \label{lem:theorem-55-triangle}
  \lean{CombinatorialRigidity.Molecular.BodyHingeFramework.theorem_55_triangle}
  \leanok
  \uses{def:rank-hypothesis, lem:rank-parallel-full}
  The three-body triangle sibling of the two-body base case
  \cref{lem:theorem-55-base}: if a body-hinge framework $F$ has three
  edges $e_1 : u$--$v$, $e_2 : v$--$w$, $e_3 : w$--$u$ forming a
  triangle whose three supporting extensors $C(p(e_1))$, $C(p(e_2))$,
  $C(p(e_3))$ are linearly independent, then $F$ is infinitesimally rigid
  on the three bodies $\{u, v, w\}$
  ($\mathrm{IsInfinitesimallyRigidOn}\,\{u,v,w\}$).
\end{lemma}
\begin{proof}
  \leanok
  \uses{lem:rank-parallel-full}
  Each hinge constraint puts a relative screw $S(u) - S(v)$, $S(v) -
  S(w)$, $S(w) - S(u)$ in the one-dimensional span of its supporting
  extensor; the three differences telescope to $0$ (the cyclic shift is a
  bijection of $\{u,v,w\}$). Linear independence of the three extensors
  --- the $m = 3$ instance of \cref{lem:rank-parallel-full}'s
  $m$-dimensional generalization --- forces each difference to vanish, so
  $S(u) = S(v) = S(w)$, and constancy on $\{u,v,w\}$ follows.
\end{proof}

\begin{lemma}[Cyclic-seed existence for the triangle base]
  \label{lem:triangle-normals}
  \lean{CombinatorialRigidity.Molecular.exists_triangle_normals}
  \leanok
  \uses{def:panel-support-extensor, lem:panel-support-extensor-independence,
        lem:extensor-independence}
  When $k \ge 1$ (so $k+2 \ge 3$), there exist three vectors $n_0, n_1, n_2 \in \R^{k+2}$
  such that (1) each cyclic pair $(n_0,n_1)$, $(n_1,n_2)$, $(n_2,n_0)$ has a nonzero
  grade-2 join ($n_i \wedge n_j \ne 0$), and (2) the cyclic supporting-extensor family
  $i \mapsto \mathrm{panelSupportExtensor}(n_{\sigma(i)}, n_{\sigma(i)+1})$ is linearly
  independent. The witness is $n_0 = e_0$, $n_1 = e_1$, $n_2 = e_2$ (standard basis
  vectors in $\R^{k+2}$): the cyclic family reduces (via the grade-2 antisymmetry at the
  reversed pair $(e_2, e_0) = -(e_0 \wedge e_2)$) to the sorted family $e_0 \wedge e_1$,
  $e_1 \wedge e_2$, $e_0 \wedge e_2$, up to unit scalars; the sorted family is a
  three-member subfamily of the $\Lambda^2$-basis indexed by distinct 2-subsets; that
  basis family is linearly independent (\cref{lem:extensor-independence}); unit scaling
  preserves linear independence; and $e_i \wedge e_j \ne 0$ for $i < j$ since each such
  join is a nonzero member of the same linearly independent basis family.
\end{lemma}
\begin{proof}
  \leanok
  \uses{lem:panel-support-extensor-independence, lem:extensor-independence}
  By \cref{lem:panel-support-extensor-independence} it suffices to show the grade-2 join
  family is linearly independent. Via \texttt{normalsJoin\_swap}, the cyclic join family
  $(e_0 \wedge e_1,\, e_1 \wedge e_2,\, e_2 \wedge e_0)$ equals (unit scalars) times the
  sorted family $(e_0 \wedge e_1,\, e_1 \wedge e_2,\, e_0 \wedge e_2)$; unit scaling
  preserves independence. Each sorted join equals the corresponding member of the
  $\Lambda^2$-basis family (\texttt{normalsJoin\_basisFun\_orderEmbOfFin}); the three
  indices $\{0,1\}$, $\{1,2\}$, $\{0,2\}$ are distinct, so the three-element subfamily
  inherits independence from the full basis family
  (\cref{lem:extensor-independence}).
\end{proof}

\begin{lemma}[Cyclic-seed existence for a general cycle]
  \label{lem:cycle-normals}
  \lean{CombinatorialRigidity.Molecular.exists_cycle_normals}
  \leanok
  \uses{def:panel-support-extensor, lem:panel-support-extensor-independence,
        lem:extensor-independence}
  For $3 \le m \le k + 2$ there exist $m$ vectors $n_0, \dots, n_{m-1} \in \R^{k+2}$
  such that (1) each cyclic pair $(n_i, n_{i+1})$ (indices mod $m$) has a nonzero
  grade-2 join $n_i \wedge n_{i+1} \ne 0$, and (2) the cyclic supporting-extensor
  family $i \mapsto \mathrm{panelSupportExtensor}(n_i, n_{i+1})$ is linearly
  independent. This is the general-$m$ generalization of the triangle seed
  \cref{lem:triangle-normals} (its $m = 3$ instance), and the \emph{shared}-normal
  sibling of the free-pair extensor existence
  \cref{lem:exists-independent-panel-extensor}: consecutive hinges of a panel cycle
  share a panel, so a panel-hinge seed assigns \emph{one} normal per body --- the
  free pairs of \cref{lem:exists-independent-panel-extensor} cannot feed it.
\end{lemma}
\begin{proof}
  \leanok
  \uses{lem:panel-support-extensor-independence, lem:extensor-independence}
  Witness: the standard basis restricted along the inclusion of index sets,
  $n_i = e_i$ (this is where $m \le k + 2$ enters). The $m$ cyclic joins then
  realize the $m$ distinct $2$-subsets $\{i, i+1\}$ ($i < m-1$) and $\{0, m-1\}$ ---
  the edges of the cycle $C_m$ embedded in the index set. Via the grade-2
  antisymmetry at the single wrap edge, the cyclic join family equals unit scalars
  times the sorted subfamily of the $\Lambda^2$-basis family indexed by these
  $2$-subsets; the edge-to-subset index map is injective for $m \ge 3$ (two cyclic
  edges coincide only if $2 = 0$ in $\Z/m$), so the subfamily inherits independence
  from the basis family (\cref{lem:extensor-independence}), and unit scaling
  preserves independence. Each join is nonzero since consecutive indices are
  distinct; the extensor conclusion is carried across
  \cref{lem:panel-support-extensor-independence}.
\end{proof}

\subsection{The framework-construction op}
\label{sec:molecular-algebraic-induction-withgraph}

Both inductive cases build a realization of $G$ from a realization of a
\emph{different} graph: Case~I from the contraction $G/E(H)$, Case~II
from the splitting-off $G_v^{ab}$. The shared, citation-free primitive
both need is the ability to keep the hinge assignment fixed while
changing the underlying multigraph --- the supporting extensors $C(p(e))$
and hinge-row blocks depend only on the hinge data, not on which edges
$G$ carries. The one fact this phase needs from the op is the
graph-monotonicity of the motion space: deleting edges (passing to a
subgraph) can only enlarge the null space $Z(G,p)$, the direction
Cases~I and~II travel toward the smaller inductive graph.

\begin{definition}[Framework on a new graph]
  \label{def:framework-with-graph}
  \lean{CombinatorialRigidity.Molecular.BodyHingeFramework.withGraph,
        CombinatorialRigidity.Molecular.PanelHingeFramework.withGraph,
        CombinatorialRigidity.Molecular.PanelHingeFramework.toBodyHinge_withGraph,
        CombinatorialRigidity.Molecular.PanelHingeFramework.withNormal,
        CombinatorialRigidity.Molecular.PanelHingeFramework.toBodyHinge_withNormal_supportExtensor_of_ne,
        CombinatorialRigidity.Molecular.PanelHingeFramework.toBodyHinge_withNormal_infinitesimalMotions_eq,
        CombinatorialRigidity.Molecular.PanelHingeFramework.toBodyHinge_withNormal_pinnedMotions_eq}
  \leanok
  \uses{def:rigidity-matrix, def:panel-hinge-framework}
  For a body-hinge framework $F = (G, p)$ and any multigraph $G'$ on the
  same bodies, $F[G']$ is the framework with underlying multigraph $G'$
  and the same hinge assignment $p$. Every supporting extensor $C(p(e))$,
  hinge-row block, and per-edge constraint is untouched; only the set of
  edges imposing a constraint changes. The panel layer carries the same
  operation on a panel-hinge framework $P$ --- swap the multigraph, keep
  the per-body panel normals and endpoint selector --- and the two commute
  with the body-hinge interpretation, $(P[G'])^\flat = (P^\flat)[G']$, so
  the graph-monotonicity and block-pin rank facts below apply verbatim to a
  panel realization re-glued onto $G$, with coplanarity preserved (the
  normals are unchanged). Case~II's $1$-extension additionally re-inserts a
  body $v$ with a \emph{chosen} panel: the per-body variant $P[v \mapsto n]$
  overrides only $v$'s panel normal, leaving the multigraph, endpoint
  selector, and every other normal fixed, so each hinge whose endpoints
  avoid $v$ keeps its supporting extensor and the inductive realization of
  the splitting-off $G_v^{ab}$ is untouched away from $v$. In particular,
  when $v$ is still unhinged (no edge of $G$ has $v$ as an endpoint, as in
  $G_v^{ab}$ before $v$'s two new edges are added), the choice of $v$'s
  panel normal changes neither the null space $Z(G,p)$ nor any
  body-$w$-pinned motion subspace, so the inductive count is preserved
  through the choice.
\end{definition}

\begin{lemma}[Edge deletion enlarges the motion space]
  \label{lem:motions-mono-of-graph-le}
  \lean{CombinatorialRigidity.Molecular.BodyHingeFramework.infinitesimalMotions_le_withGraph_of_le,
        CombinatorialRigidity.Molecular.BodyHingeFramework.finrank_infinitesimalMotions_le_of_graph_le}
  \leanok
  \uses{def:framework-with-graph}
  If $G' \le G$ then $Z(G,p) \subseteq Z(G', p)$, equivalently
  $\dim Z(G,p) \le \dim Z(G', p)$ --- i.e.
  $\operatorname{rank} R(G', p) \le \operatorname{rank} R(G, p)$. A
  supergraph imposes more hinge constraints, so a motion of $G$ is a
  fortiori a motion of any subgraph $G'$. This is the
  ``re-adding edges only grows the rank'' monotonicity that lifts a
  realization of a minimal $k$-dof spanning subgraph to one of the whole
  multigraph (the step \cref{prop:rigidity-matrix-prop11} uses), and the
  companion to the fixed-graph span-monotonicity Lemma~5.2.
\end{lemma}
\begin{proof}
  \leanok
  \uses{def:framework-with-graph}
  Each link of $G'$ is a link of $G$ (\texttt{Graph.IsLink.mono}), and
  the matching supporting extensors agree, so every infinitesimal motion
  of $G$ meets every constraint of $G'$: $Z(G,p) \subseteq Z(G', p)$.
  Taking dimensions (\texttt{Submodule.finrank\_mono}) gives the rank
  form.
\end{proof}


\subsection{Case I: a proper rigid subgraph}
\label{sec:molecular-algebraic-induction-caseI}

This subsection proves Katoh--Tanigawa's Case~I: a minimal $k$-dof-graph
with a proper rigid subgraph $H$ realizes its target rank by contracting
$H$ to a single pinned body and gluing the rigid block to the
contraction's own inductive realization (KT eq.~(6.3)). Two pieces of
shared infrastructure come first: the panel-cycle realization (KT
Lemma~5.4), used later in Case~III, and the block pin --- the
framework-side model of contracting a rigid block to one body. The Case~I
accounting itself (\cref{lem:case-I}) is genericity-free; genericity
enters only through KT Claim~6.4 (\cref{lem:claim-6-4}), needed to realize
both legs of the block-triangular splice at a single common placement.
The realization lemmas (\cref{lem:case-I-realization}, \cref{lem:case-I-realization-nonsimple},
\cref{lem:case-I-realization-h65}, \cref{lem:case-I-dispatch}) assemble
this into KT Lemmas~6.2, 6.3, and 6.5.

Katoh--Tanigawa's Lemma~5.4 realizes a cycle graph directly: for
$3 \le m \le D$, a cycle of $m$ bodies whose consecutive panels have
linearly independent supporting extensors is infinitesimally rigid at the
full rank $D(m-1)$. The argument has four pieces: the two-body base case;
the dimension bound $m \le D$ forced by independence in the
$D$-dimensional screw space; the telescoping argument generalizing the
base case to any cycle length; and the existence of $m$ independent
panels for $3 \le m \le D$, an exterior-algebra statement on the
grade-$2$ joins of the panel normals that a basis choice on
$\Lambda^2 \R^{k+2}$ settles.

\begin{lemma}[Cycle realization, short-cycle base]
  \label{lem:cycle-realization-base}
  \lean{CombinatorialRigidity.Molecular.PanelHingeFramework.toBodyHinge_rankHypothesis_zero}
  \leanok
  \uses{def:panel-hinge-framework, def:rank-hypothesis, lem:theorem-55-base}
  A panel-hinge framework $P$ with two edges $e_1, e_2$ joining two
  distinct bodies $u \ne v$, whose panel support extensors
  $C(p(e_1)), C(p(e_2))$ are linearly independent, has a body-hinge
  interpretation infinitesimally rigid on $\{u,v\} = V(G)$
  (\cref{def:rank-hypothesis}), the full rank $D(2-1) = D$.
\end{lemma}
\begin{proof}
  \leanok
  \uses{lem:theorem-55-base}
  The body-hinge interpretation carries exactly the two independent
  supporting extensors as its hinge data, so the parallel-hinges-full base
  case (\cref{lem:theorem-55-base}) applies verbatim.
\end{proof}

\begin{lemma}[Rigid cycle dimension bound]
  \label{lem:cycle-realization-dim-bound}
  \lean{CombinatorialRigidity.Molecular.card_le_screwDim_of_linearIndependent,
        CombinatorialRigidity.Molecular.PanelHingeFramework.card_le_screwDim_of_supportExtensor_linearIndependent}
  \leanok
  \uses{def:panel-hinge-framework, def:rigidity-matrix}
  If the supporting extensors of $m$ edges of a panel-hinge framework are
  linearly independent in the screw space $\Lambda^k \R^{k+2}$, then
  $m \le D = \binom{k+2}{2}$.
\end{lemma}
\begin{proof}
  \leanok
  A linearly independent family in a $D$-dimensional space has at most $D$
  members, and $\dim \Lambda^k \R^{k+2} = D$ (\cref{def:rigidity-matrix}).
\end{proof}

\begin{lemma}[Panel independence reduces to grade-$2$ join independence]
  \label{lem:panel-support-extensor-independence}
  \lean{CombinatorialRigidity.Molecular.panelSupportExtensor_linearIndependent_iff}
  \leanok
  \uses{def:panel-support-extensor, def:meet-complement-iso}
  A family of $m$ panel support extensors
  $i \mapsto C(p(e_i)) = C(n_{u_i}, n_{v_i})$ is linearly independent in
  the screw space $\Lambda^k \R^{k+2}$ iff the family of grade-$2$ joins
  $i \mapsto n_{u_i} \wedge n_{v_i}$ is independent in
  $\Lambda^2 \R^{k+2}$.
\end{lemma}
\begin{proof}
  \leanok
  The complement iso $\Lambda^2 \R^{k+2} \simeq \Lambda^k \R^{k+2}$
  (\cref{def:meet-complement-iso}) is a linear equivalence, and the panel
  support extensor is its image of the grade-$2$ join by definition
  (\cref{def:panel-support-extensor}); a linear equivalence both
  preserves and reflects linear independence.
\end{proof}

\begin{lemma}[Existence of independent panels for $m \le D$]
  \label{lem:exists-independent-panel-extensor}
  \lean{CombinatorialRigidity.Molecular.exists_independent_panelSupportExtensor,
        CombinatorialRigidity.Molecular.exists_independent_normalsJoin}
  \leanok
  \uses{def:panel-support-extensor, lem:panel-support-extensor-independence,
        lem:extensor-independence}
  For every $m \le D = \binom{k+2}{2}$ there are $m$ pairs of panel
  normals $(n_{u_i}, n_{v_i})$ whose supporting extensors are linearly
  independent in the screw space $\Lambda^k \R^{k+2}$.
\end{lemma}
\begin{proof}
  \leanok
  By \cref{lem:panel-support-extensor-independence} it suffices to
  produce $m$ independent grade-$2$ joins. The index set of $2$-element
  subsets of $\{1, \dots, k+2\}$ has cardinality $\binom{k+2}{2} = D \ge
  m$; picking $m$ distinct $2$-subsets and taking the corresponding
  pairs of standard basis vectors gives $m$ joins forming a subfamily of
  the basis-indexed exterior-power family of $\Lambda^2 \R^{k+2}$, which
  is linearly independent by the extensor-independence Lemma~2.1
  (\cref{lem:extensor-independence}); a subfamily of an independent
  family is independent.
\end{proof}

\begin{lemma}[Rigidity of a cycle with independent extensors]
  \label{lem:cycle-realization-rigid}
  \lean{CombinatorialRigidity.Molecular.PanelHingeFramework.toBodyHinge_rankHypothesis_zero_cycle,
        CombinatorialRigidity.Molecular.BodyHingeFramework.rankHypothesis_zero_of_cycle,
        CombinatorialRigidity.Molecular.BodyHingeFramework.isTrivialMotion_of_isInfinitesimalMotion_cycle,
        CombinatorialRigidity.Molecular.BodyHingeFramework.theorem_55_cycle}
  \leanok
  \uses{def:panel-hinge-framework, def:rank-hypothesis, lem:cycle-realization-base}
  A panel-hinge framework on the cycle of $m$ bodies $0, 1, \dots, m-1$
  --- the $i$-th hinge joining body $i$ to body $i+1$ (cyclically) ---
  whose $m$ supporting extensors are linearly independent in the screw
  space $\Lambda^k \R^{k+2}$ is infinitesimally rigid, realizing the rank
  hypothesis at $k' = 0$ (\cref{def:rank-hypothesis}). The same conclusion
  holds when the cycle is indexed by an arbitrary map
  $\mathrm{vtx} \colon \Z/m \to V$, with no injectivity required: rigidity
  is then asserted on the bodies the cycle actually visits,
  $\operatorname{range}(\mathrm{vtx})$.
\end{lemma}
\begin{proof}
  \leanok
  \uses{lem:rank-parallel-full}
  Each hinge constraint puts the relative screw $S(i) - S(i+1)$ of an
  infinitesimal motion $S$ in the one-dimensional span of its supporting
  extensor, and the $m$ differences telescope around the cycle to
  $\sum_i (S(i) - S(i+1)) = 0$. Independence of the $m$ extensors then
  forces every difference to vanish --- the $m$-edge generalization of
  the parallel-hinges-full Lemma~5.3 (\cref{lem:rank-parallel-full}, the
  $m=2$ case) --- so $S$ is constant around the connected cycle, hence a
  trivial motion, and \cref{def:rank-hypothesis} gives the rank hypothesis
  at $k'=0$.
\end{proof}

\begin{lemma}[Cycle realization; KT Lemma 5.4]
  \label{lem:cycle-realization}
  \lean{CombinatorialRigidity.Molecular.PanelHingeFramework.cycle_realization}
  \leanok
  \uses{def:rigidity-matrix, def:rank-hypothesis, def:cycle-data, def:k-dof,
        lem:cycle-realization-base, lem:cycle-realization-dim-bound,
        lem:cycle-realization-rigid, lem:cycle-normals,
        lem:case-III-claim612-line-in-panel-union, lem:genericity-device}
  A minimal $0$-dof-graph $G$ that \emph{is} a cycle graph
  (\cref{def:cycle-data}) on $m$ vertices, $3 \le m \le D$, has a generic
  full-rank panel-hinge realization $(G,p)$ (\cref{def:rank-hypothesis}):
  general-position, algebraically independent panel normals attaining the
  full rank $D(|V|-1)$.
\end{lemma}
\begin{proof}
  \leanok
  \uses{lem:cycle-realization-rigid,
        lem:cycle-normals, lem:case-III-claim612-line-in-panel-union}
  The cyclic shared-normal seed (\cref{lem:cycle-normals}) supplies $m$
  panel normals --- one per body, consecutive hinges sharing a panel ---
  whose $m$ cyclic supporting extensors are linearly independent and
  whose cyclic grade-$2$ joins are nonzero ($m \le d = k+1$, the
  dimension chain from $D = \binom{d+1}{2} = \binom{k+2}{2}$). Assign
  these normals to the cycle vertices and orient each edge by the
  graph's canonical endpoint choice: every cycle edge's supporting
  extensor is then $\pm$ a member of the seed's independent cyclic
  family, so the edge-extensor family is independent and nonvanishing
  edge-by-edge (every edge of $G$ \emph{is} a cycle edge). The
  telescoping rigidity (\cref{lem:cycle-realization-rigid}, in its
  arbitrary-body-type form) makes this realization infinitesimally rigid
  on the cycle bodies, i.e.\ on $V(G)$, and the bare-to-generic upgrade
  (\cref{lem:case-III-claim612-line-in-panel-union}) re-realizes it
  generically at full rank, exactly as for the triangle base
  (\cref{lem:triangle-realization}).
\end{proof}

\begin{fmlnote}
  KT states Lemma~5.4 for the full range $3 \le m \le D$; the formalized
  statement covers $m \le d$ (with $D = \binom{d+1}{2}$), the range the
  chain-extraction dichotomy of \cref{sec:molecular-algebraic-induction-caseIII}
  actually consumes --- the remaining band $d < m \le D$ has no consumer in
  this induction. Only the ``independent $\Rightarrow$ rigid'' direction of
  the Crapo--Whiteley equivalence is used (KT realizes the cycle \emph{at} an
  independent hinge family); the projective assembly of the realization on a
  concrete cycle graph, which KT cites rather than reproves, is also carried
  out here. This is Katoh--Tanigawa's Lemma~5.4; its geometric content is due
  to Crapo--Whiteley~\cite[Proposition~3.4]{crapoWhiteley1982} (the $d=3$
  case) and Whiteley~\cite[Proposition~3]{whiteley1999}, the argument
  extending to general $d$ without modification.
\end{fmlnote}

The block pin is the framework-side model of contracting a set of bodies
to a single pinned one, the move Case~I makes on a rigid subgraph.

\begin{definition}[Block-pinning a set of bodies]
  \label{def:pinned-motions-on}
  \lean{CombinatorialRigidity.Molecular.BodyHingeFramework.pinnedMotionsOn,
        CombinatorialRigidity.Molecular.BodyHingeFramework.pinnedMotionsOn_singleton,
        CombinatorialRigidity.Molecular.BodyHingeFramework.pinnedMotionsOn_eq_iInf}
  \leanok
  \uses{def:rigidity-matrix, lem:rank-delete-vertex}
  For a body-hinge framework $F = (G,p)$ and a set of bodies
  $s \subseteq V$, the \emph{block pin} $Z_s(G,p)$ is the subspace of
  infinitesimal motions vanishing on every body of $s$:
  $\{S \in Z(G,p) : S|_s = 0\}$.
\end{definition}

At a singleton $s = \{v\}$ the block pin recovers the pin-a-body subspace
$\operatorname{pinned}(v)$ of \cref{lem:rank-delete-vertex}, and for a
nonempty block it equals the intersection of the single-body pins,
$Z_s(G,p) = \bigsqcap_{v \in s} \operatorname{pinned}(v)$. Case~I applies
the block pin to a rigid subgraph $H$: pinning all of $V(H)$ to the zero
screw is the algebraic effect of contracting $H$ to a single body.

\begin{lemma}[Edge deletion enlarges the block pin]
  \label{lem:pinned-motions-on-mono}
  \lean{CombinatorialRigidity.Molecular.BodyHingeFramework.pinnedMotionsOn_le_withGraph_of_le,
        CombinatorialRigidity.Molecular.BodyHingeFramework.finrank_pinnedMotionsOn_le_of_graph_le}
  \leanok
  \uses{def:pinned-motions-on, def:framework-with-graph, lem:motions-mono-of-graph-le}
  Replacing $G$ by any subgraph $G' \le G$ (keeping the hinge data,
  \cref{def:framework-with-graph}) can only grow the block pin:
  $Z_s(G,p) \subseteq Z_s(G',p)$, hence
  $\dim Z_s(G,p) \le \dim Z_s(G',p)$.
\end{lemma}
\begin{proof}
  \leanok
  \uses{lem:motions-mono-of-graph-le}
  The block pin $Z_s(G,p)$ is $Z(G,p)$ intersected with the
  graph-independent vanishing conditions $S|_s = 0$, so the inclusion is
  the motion-space monotonicity (\cref{lem:motions-mono-of-graph-le})
  applied to the first factor; the dimension bound follows.
\end{proof}

This is the direction Case~I's block-triangular gluing travels: placing
the contraction realization on the smaller graph $G/E(H)$ and re-adding
the edges $E(H)$ only grows the block-pinned rank, the slack in
\cref{lem:pinned-motions-on-rank-bound} filled by the contraction's own
rank.

\begin{lemma}[Block-pinning drops at least the $D$ trivial-motion dimensions]
  \label{lem:pinned-motions-on-rank-bound}
  \lean{CombinatorialRigidity.Molecular.BodyHingeFramework.screwDim_add_finrank_pinnedMotionsOn_le}
  \leanok
  \uses{def:pinned-motions-on, lem:rank-delete-vertex, lem:trivial-motions-rank-bound}
  For a nonempty block of bodies $s \subseteq V$,
  \[
    D + \dim Z_s(G,p) \;\le\; \dim Z(G,p).
  \]
\end{lemma}
\begin{proof}
  \leanok
  \uses{lem:rank-delete-vertex}
  The block pin $Z_s(G,p)$ sits inside the single-body pin
  $\operatorname{pinned}(v)$ of any $v \in s$, which meets the
  $D$-dimensional trivial motions only at $0$
  (\cref{lem:rank-delete-vertex}); so the block pin and the trivial
  motions are disjoint submodules of $Z(G,p)$, and their dimensions add
  to at most $\dim Z(G,p)$.
\end{proof}

This is the lower-bound half of Case~I's block-triangular accounting: a
single-body pin is an exact $+D$ direct-sum split (\cref{lem:rank-delete-vertex}),
whereas pinning a whole rigid block $H$ collapses it to one effective
body and the residual $D(|V(H)|-1)$ constraints make the bound an
inequality --- the contraction's own rank, supplied by the induction
hypothesis, fills the slack to equality. The next lemma records the two
extreme cases of that slack.

\begin{lemma}[The block pin vanishes exactly where the framework is rigid]
  \label{lem:pinned-motions-on-rigid-iff}
  \lean{CombinatorialRigidity.Molecular.BodyHingeFramework.pinnedMotionsOn_eq_bot_of_isInfinitesimallyRigid,
        CombinatorialRigidity.Molecular.BodyHingeFramework.finrank_pinnedMotionsOn_eq_zero_of_isInfinitesimallyRigid,
        CombinatorialRigidity.Molecular.BodyHingeFramework.isInfinitesimallyRigid_of_block_of_pinnedMotionsOn_eq_bot}
  \leanok
  \uses{def:pinned-motions-on}
  If $F$ is infinitesimally rigid, then block-pinning any nonempty body
  set $s$ leaves nothing: $Z_s(G,p) = 0$. Conversely, if every
  infinitesimal motion of $F$ is constant on a nonempty block $s$ and the
  block pin vanishes, $Z_s(G,p) = 0$, then $F$ is infinitesimally rigid.
\end{lemma}
\begin{proof}
  \leanok
  A block-pinned motion is an infinitesimal motion, so rigidity makes it
  trivial (constant); vanishing at one body of $s$ then forces the
  constant to be $0$. Conversely, given a motion $S$ and a body
  $v \in s$, the constant assignment equal to $S(v)$ everywhere is a
  trivial motion; their difference vanishes on all of $s$ by the
  block-constancy hypothesis, hence lies in the (vanishing) block pin,
  forcing $S$ to equal that constant.
\end{proof}

\begin{lemma}[Case I: rigid-subgraph contraction realizes the rank; KT Lemmas 6.2, 6.3, 6.5]
  \label{lem:case-I}
  \lean{CombinatorialRigidity.Molecular.BodyHingeFramework.isInfinitesimallyRigidOn_iff_pinnedMotionsOn_le,
        CombinatorialRigidity.Molecular.PanelHingeFramework.toBodyHinge_rankHypothesis_iff_finrank_pinnedMotionsOn,
        CombinatorialRigidity.Molecular.BodyHingeFramework.rankHypothesis_iff_finrank_pinnedMotionsOn}
  \leanok
  \uses{thm:minimal-kdof-reduction, lem:contraction-minimality,
        lem:rank-delete-vertex, def:pinned-motions-on,
        lem:pinned-motions-on-rank-bound, lem:genericity-device}
  Let $G$ be a minimal $k$-dof-graph with a proper rigid subgraph $H$ on
  the (nonempty) body set $s = V(H)$. Contracting $H$ to a single pinned
  body is the block pin $Z_s(G,p)$ (\cref{def:pinned-motions-on}), the
  framework-side carrier of the contraction $G/E(H)$ (a smaller minimal
  $k$-dof-graph by \cref{lem:contraction-minimality}, supplied the
  induction hypothesis). Then $(G,p)$ realizes the $V(G)$-relative target
  rank at $k'$ (\cref{def:rank-hypothesis}) \emph{iff} the rigid block $H$
  contributes its full $D(|V(H)|-1)$ rows and the residual contraction
  contributes the rest --- the block-triangular rank addition of
  KT~(6.3): $\operatorname{rank} R(G,p) = D(|V(H)|-1) +
  \operatorname{rank} R(G/E(H), p_2)$, glued via the pin-a-body
  Lemma~5.1 (\cref{lem:rank-delete-vertex}).
\end{lemma}
\begin{proof}
  \leanok
  \uses{lem:rank-delete-vertex, lem:pinned-motions-on-rank-bound,
        lem:genericity-device}
  Pin the rigid block $H$ and place the contraction $G/E(H)$ at its
  inductive rank. The block pin and the $D$-dimensional trivial motions
  are disjoint, giving the genericity-free lower bound
  (\cref{lem:pinned-motions-on-rank-bound}). The combined rigidity matrix
  is block-triangular --- the $H$-rows occupy only the $V(H)$ columns,
  the contraction rows the rest --- so at a generic point the ranks add
  and the reverse bound holds (the entries are polynomials in
  algebraically independent panel coordinates, so a generic point attains
  the maximal rank, KT's Claim~6.4, \cref{lem:genericity-device}). The
  two bounds together pin the parent rank to $D(|V(G)|-1) - k$. This
  accounting depends only on $s, t \subseteq V(G)$, not on the ambient
  body type, so it composes through the vertex-reducing induction; the
  companion nullity form (relating $\dim Z(G,p)$ to the ambient body
  type) remains available for the deficiency and Proposition~1.1 path.
\end{proof}

\begin{lemma}[Panel-coordinate row coordinatization; KT Claim 6.4]
  \label{lem:rows-polynomial-in-normals}
  \lean{CombinatorialRigidity.Molecular.PanelHingeFramework.exists_good_realization_ofParam}
  \leanok
  \uses{def:panel-support-extensor, def:rigidity-matrix, lem:genericity-device}
  Regard a panel-hinge realization of a fixed graph $G$ as a point of
  panel-coordinate space: an assignment $q \colon V_\alpha \times \{1,
  \dots, k+2\} \to \R$ of a normal $n_v = q(v, -) \in \R^{k+2}$ to each
  body. Then the rigidity rows of $R(G,q)$ form a finite family of
  functionals on $\R^{D|V_\alpha|}$ whose coordinates, in the standard
  basis, are degree-$2$ polynomials in the entries of $q$ --- exactly the
  input shape the genericity device (\cref{lem:genericity-device}) needs
  to be applied to a \emph{varying} panel realization, rather than a
  constant one.
\end{lemma}
\begin{proof}
  \leanok
  A coordinate of the grade-$2$ join $n_u \wedge n_v$ against the
  standard dual basis of $\Lambda^2 \R^{k+2}$ is the $2 \times 2$ minor
  $n_u(i)\,n_v(j) - n_u(j)\,n_v(i)$, a determinant of the two normals'
  coordinates and hence bilinear, so it is a degree-$2$ polynomial in the
  panel coordinates $q$. The supporting extensor $C(p(e)) =
  \mathrm{complementIso}(n_u \wedge n_v)$ (\cref{def:panel-support-extensor})
  is a fixed linear image of the join, so its coordinates stay degree-$2$.
  Each rigidity row spans the annihilator of its edge's supporting
  extensor in $(\R^{k+2})^*$, and its coordinates are in turn a fixed
  linear combination of the extensor's coordinates at the two endpoints,
  so they too are degree-$2$ in $q$. Applying the genericity device to
  this family, with the witnessed independent subfamily supplied by the
  general-position seed (\cref{lem:exists-independent-panel-extensor}),
  produces a normal assignment attaining the codimension bound
  $\#s + \dim Z(G,q) \le D|V_\alpha|$ at the same corank.
\end{proof}

\begin{lemma}[Case I splice glue; KT eq.\ (6.3)]
  \label{lem:case-I-splice-seed}
  \lean{CombinatorialRigidity.Molecular.BodyHingeFramework.isInfinitesimallyRigidOn_of_splice}
  \leanok
  \uses{def:framework-with-graph, def:rank-hypothesis, lem:case-I}
  Let $H$ (on $s_H = V(H)$) be a proper rigid subgraph and $G/E(H)$ (on
  $s_c$) its contraction, both subgraphs of $G$ sharing the contracted
  body $c \in s_H \cap s_c$ and covering $V(G) \subseteq s_H \cup s_c$.
  If $F$ restricted to $H$ (\cref{def:framework-with-graph}) is
  infinitesimally rigid on $s_H$, and $F$ restricted to $G/E(H)$ is
  infinitesimally rigid on $s_c$, then $F$ is infinitesimally rigid on
  $V(G)$.
\end{lemma}
\begin{proof}
  \leanok
  \uses{lem:motions-mono-of-graph-le}
  Each leg transports to the parent by the monotonicity that re-adding
  edges only shrinks the motion space (\cref{lem:motions-mono-of-graph-le},
  since $H, G/E(H) \le G$), so a parent motion is already constant on
  $s_H$ and on $s_c$. The two constants agree at the shared contracted
  body $c$, so every motion is constant across $s_H \cup s_c \supseteq
  V(G)$.
\end{proof}

This is KT's block-triangular splice (eq.~(6.3)) read as two overlapping
rigid pieces gluing along a shared body, and it is genericity-free: the
genericity device (\cref{lem:genericity-device}) enters only at the
realization lemma \cref{lem:case-I-realization}, where the two legs must be
realized at a single generic placement --- the intersection of the two
legs' Zariski-open rigid loci. The hypotheses here are the satisfiable
inductive facts, relative rigidity of each piece on a common placement
$F$, not the parent rank they conclude.

\begin{lemma}[Graph-level minimality of the rigid-subgraph contraction]
  \label{lem:rigidContract-isMinimalKDof}
  \lean{Graph.rigidContract_isMinimalKDof}
  \leanok
  \uses{lem:contraction-minimality, def:rigid-contraction, lem:reduction-step}
  Contracting a proper rigid subgraph $H$ of a minimal $k$-dof-graph $G$
  again yields a minimal $k$-dof-graph: $(G/E(H))$ is minimal $k$-dof,
  with strictly fewer vertices (\cref{lem:reduction-step}).
\end{lemma}
\begin{proof}
  \leanok
  \uses{lem:contraction-minimality, thm:def-eq-corank, lem:two-edge-conn}
  The reduction theorem (\cref{thm:minimal-kdof-reduction}) supplies only
  the existence of a rigid subgraph together with the induction
  hypothesis for the reduction, leaving the contraction's own minimality
  to be established separately; that is the content here. The matroid
  contraction $M(\tilde G)/E(\tilde H)$ is already a minimal $k$-dof
  matroid (\cref{lem:contraction-minimality}), and the matroid of the
  \emph{graph} contraction $G/E(H)$ equals that matroid contraction ---
  the graph deletes the edges $E(H)$ and relabels vertices, while the
  matroid contracts the fibers $E(\tilde H)$; reconciling the two needs
  $\tilde H$ connected on $V(H)$, which follows from $H$ being $0$-dof
  (\cref{lem:two-edge-conn}). With that identification the base/fiber
  condition transports directly (an edge of $G/E(H)$ is a surviving
  $G$-edge), and the deficiency transports via the def-equals-corank
  bridge (\cref{thm:def-eq-corank}) against the conserved matroid rank and
  the collapse vertex-count $|V(G/E(H))| = |V(G)|-|V(H)|+1$.
\end{proof}

\begin{lemma}[KT Claim 6.4: the surviving block attains its rank at a generic placement (eq.\ (6.9))]
  \label{lem:claim-6-4}
  \lean{CombinatorialRigidity.Molecular.PanelHingeFramework.rigidContract_exterior_rank_transport_htransport}
  \leanok
  \uses{def:rigidity-matrix, def:rigid-contraction, lem:rank-delete-vertex,
        lem:contraction-minimality}
  Let $G$ be a simple minimal $0$-dof-graph with a proper rigid subgraph
  $H$ on $s = V(H)$, and write $s_c = (V(G) \setminus V(H)) \cup \{r\}$
  for the surviving body set with representative $r$. Given the
  contraction $G/E(H)$ realized at its inductive rank, the surviving-edge
  block of the rigidity matrix --- restricted to the edges $E(G)
  \setminus E(H)$ and to the columns outside $V(H)$ --- attains rank
  $D(|s_c|-1)$ at a generic placement: at a Zariski-open locus of panel
  coordinates, the surviving rows are linearly independent after
  projecting away the columns of $V(H)$ (KT's bottom-right block,
  eq.~(6.9); the row-count companion of the pin-a-body Lemma~5.1,
  \cref{lem:rank-delete-vertex}).
\end{lemma}
\begin{proof}
  \leanok
  \uses{lem:rank-delete-vertex, lem:contraction-minimality}
  The argument transports the contraction's rank across the collapse map
  $f$ sending every body of $V(H)$ to $r$ and fixing the rest, from a
  generic realization of the abstract relabelled graph $G/E(H)$ to the
  surviving-edge block of the parent.

  First, the contraction's generic realization is transported to the
  endpoint selector recording the links of $G/E(H)$ through $f$: since
  both selectors record the same links, they agree up to swap and the two
  realizations have the same rigidity and motion space. This gives a
  free-normal panel framework on $G/E(H)$ in general position, rigid on
  its whole vertex set $s_c$.

  Second --- the genuine content of the claim --- the exterior-column
  projection that drops exactly the $V(H)$-columns loses no rank on this
  framework's rigidity-row span. Because the framework is rigid on $s_c$
  and $s_c \cap V(H) = \{r\}$, the trivial motions and the projection's
  kernel together span everything, so the projected rows stay
  independent, of size at least $D(|s_c|-1)$: KT's bottom-right block,
  the pin-a-body content (\cref{lem:rank-delete-vertex}) at the row
  level.

  Third, the witness placement is KT's degenerate seed (eq.~(6.7)): the
  collapsed normal field $n \circ f$ on $G \setminus E(H)$. A per-edge
  identity --- both framings read the same support extensor, the
  projection reconciling their differing endpoints --- carries the
  projected independence of the relabelled contraction's rows back to
  the projected independence of the parent's surviving-edge rows at this
  seed.

  This discharges what KT \S5.1/Claim~6.4 states under the simplifying
  assumption that the panel coordinates are algebraically independent
  over $\Q$, realizing that genericity as the contraction's inductive
  general-position realization, the transport as the endpoint-selector
  agreement, and the rank preservation as the zero-rank-loss of the
  exterior-column projection on a rigid block.
\end{proof}

\begin{lemma}[Case I realization: simple-contraction case, all deficiencies; KT Lemma~6.3]
  \label{lem:case-I-realization}
  \lean{CombinatorialRigidity.Molecular.PanelHingeFramework.case_I_realization_all_k_gen}
  \leanok
  \uses{def:rank-hypothesis, lem:case-I, lem:rigidContract-isMinimalKDof,
        lem:rank-polynomial-IH-relabel-proj, lem:genericity-device,
        lem:rows-polynomial-in-normals, lem:isInfRigidOn-of-relative-count}
  Let $G$ be a simple minimal $k$-dof-graph with $|V(G)| \ge 3$, a proper
  rigid subgraph $H$ on $s = V(H)$ with representative body $r$, such
  that the contraction $G/E(H)$ is itself simple (KT Lemma~6.3's case
  hypothesis; the non-simple-contraction case is
  \cref{lem:case-I-realization-h65}). If the induction hypothesis
  supplies generic realizations of both $H$ (at deficiency $0$, since $H$
  is rigid) and the contraction $G/E(H)$ (at deficiency $k$, minimal
  $k$-dof by \cref{lem:rigidContract-isMinimalKDof}), then $G$ carries a
  generic realization at rank $D(|V(G)|-1) - k$ (\cref{def:rank-hypothesis}),
  discharging the Case~I contraction premise of \cref{thm:theorem-55} in
  this case.
\end{lemma}
\begin{proof}
  \leanok
  \uses{lem:rows-polynomial-in-normals, lem:rank-polynomial-IH-relabel-proj,
        lem:genericity-device, lem:case-I, lem:rigidContract-isMinimalKDof,
        lem:isInfRigidOn-of-relative-count}
  Work on a single framework $F$ at a general-position seed $q_0$, KT's
  block-triangular rank addition (eq.~(6.3)). The rigid-block rows
  ($E(H)$) are independent and number $D(|V(H)|-1)$, from the $H$-leg's
  inductive generic realization at deficiency $0$, fed through the row
  coordinatization of \cref{lem:rows-polynomial-in-normals}. The
  surviving-edge rows ($E(G) \setminus E(H)$, on $s_c = (V(G) \setminus
  V(H)) \cup \{r\}$) number $D(|s_c|-1) - k$ and, after the
  exterior-column projection killing the $V(H)$ columns, are independent
  (\cref{lem:rank-polynomial-IH-relabel-proj}, the deficiency-aware
  surviving-block rank polynomial applied to the contraction's inductive
  realization). The rigid-block rows vanish under this projection (KT's
  top-right zero block), so the two row families are disjoint and their
  union gives $D(|V(H)|-1) + D(|s_c|-1) - k = D(|V(G)|-1) - k$
  independent rows (using $|s_H \cap s_c| = 1$). The genericity device
  (\cref{lem:genericity-device}) then reads rigidity on $V(G)$ off this
  row count via the relative-count adapter
  (\cref{lem:isInfRigidOn-of-relative-count}), at $q_0$ itself; since
  $q_0$ is general position, the realization is generic.
\end{proof}

\begin{lemma}[Case I realization: non-simple $G$; KT Lemma~6.2]
  \label{lem:case-I-realization-nonsimple}
  \lean{CombinatorialRigidity.Molecular.case_I_realization_nonsimple}
  \leanok
  \uses{def:genuine-hinge-realization, lem:case-I, lem:two-vertex-zero-dof,
        lem:rigidContract-isMinimalKDof}
  Let $G$ be a minimal $k$-dof-graph with $|V(G)| \ge 3$ that is
  \emph{not} simple: it carries a parallel pair, edges $e \ne f$ both
  linking vertices $a \ne b$. Given the induction hypothesis for every
  minimal $k'$-dof-graph with fewer vertices, $G$ has a nondegenerate-hinge
  panel realization at rank $D(|V(G)|-1) - k$
  (\cref{def:genuine-hinge-realization}).
\end{lemma}
\begin{proof}
  \leanok
  \uses{lem:extensor-pair-in-panel, lem:theorem-55-base,
        lem:trivial-motions-rank-bound, lem:rigidityRows-splice-rank-add,
        lem:extProj-preserves-rank-of-inter}
  The parallel pair $\{e,f\}$ spans a proper rigid subgraph
  $H' = G[\{a,b\}] \restriction \{e,f\}$, $0$-dof by
  \cref{lem:two-vertex-zero-dof} and proper since $|V(G)| \ge 3$.
  Contract to $G' = G/H'$ and apply the induction hypothesis to obtain a
  panel-hinge framework $F_c$ on $G'$; install coincident panels at $a$
  and $b$ by pulling back $F_c$'s normal at the contracted body, and
  choose two linearly independent extensors $C_e, C_f$ in that shared
  panel (\cref{lem:extensor-pair-in-panel}). The two-body base case
  (\cref{lem:theorem-55-base}) makes the $H'$-block rigid at rank $D$;
  projecting away the $H'$-columns, the $H'$-rows vanish while $F_c$'s
  rows inject (\cref{lem:extProj-preserves-rank-of-inter}), so the two
  blocks are jointly independent and $D + \operatorname{rank}(F_c) \le
  \operatorname{rank}(F)$ (\cref{lem:rigidityRows-splice-rank-add}). The
  upper bound $\operatorname{rank}(F) \le D(|V(G)|-1)-k$ follows from the
  codimension bound (\cref{lem:trivial-motions-rank-bound}); the
  arithmetic $D + (D(|V(G)|-2)-k) = D(|V(G)|-1)-k$ closes rank equality.
\end{proof}

\begin{lemma}[Case I realization: removing a degree-two vertex; KT Lemma~6.5, Claim~6.6]
  \label{lem:case-I-realization-h65}
  \lean{CombinatorialRigidity.Molecular.PanelHingeFramework.case_I_realization_h65_gen}
  \leanok
  \uses{def:rank-hypothesis, lem:genericity-device}
  Let $G$ be a simple minimal $0$-dof-graph with $|V(G)| \ge 3$ in which
  every proper rigid subgraph has a non-simple contraction. Given the
  induction hypothesis for every smaller minimal $0$-dof-graph, $G$
  carries a generic full-rank realization (\cref{def:rank-hypothesis}).
\end{lemma}
\begin{proof}
  \leanok
  \uses{lem:genericity-device}
  By KT Claim~6.6, a vertex-inclusionwise maximal proper rigid subgraph
  $G'$ of $G$ has a non-simple contraction only because some vertex
  $v \notin V(G')$ has degree two, with two neighbours $a, b \in V(G')$
  and $G = G' + v$; then $G - v$ is a smaller simple minimal $0$-dof-graph.
  The induction hypothesis gives a generic full-rank realization of
  $G - v$ with algebraically independent seed $q$. Re-attach $v$ at the
  same seed: keep every normal and endpoint from $G-v$'s realization,
  give $v$ its own normal $q(v)$, and direct the two new edges $va, vb$
  from $v$. Off $\{va, vb\}$ the framework agrees with the inductive
  realization, so its $D(|V(G-v)|-1)$ independent rigidity rows transport
  unchanged, vanishing on $v$'s screw column. The two new rows at $va,
  vb$ span the full $D$-dimensional block through $v$'s column: their
  supporting extensors $C(q(v), q(a))$ and $C(q(v), q(b))$ are linearly
  independent because the triple $(q(v), q(a), q(b))$ is --- $q$ being
  algebraically independent and $G$ having only the three vertices $v,
  a, b$ forces this by a direct computation, with no perturbation
  argument needed. The two row blocks are disjoint (one vanishes on,
  the other is pinned through, $v$'s column) and together give $D +
  D(|V(G-v)|-1) = D(|V(G)|-1)$ independent rows, so $G$ is
  infinitesimally rigid on $V(G)$ at this seed; the genericity device
  (\cref{lem:genericity-device}) upgrades this fixed-seed rigidity to a
  generic realization.
\end{proof}

\begin{lemma}[Case I dispatch: KT Lemmas 6.2, 6.3, 6.5]
  \label{lem:case-I-dispatch}
  \lean{CombinatorialRigidity.Molecular.case_I_dispatch_gen}
  \leanok
  \uses{def:genuine-hinge-realization, def:rank-hypothesis,
        lem:case-I-realization-nonsimple, lem:case-I-realization,
        lem:case-I-realization-h65}
  Let $G$ be a minimal $0$-dof-graph with $|V(G)| \ge 3$ carrying a
  proper rigid subgraph. Then $G$ has a nondegenerate-hinge panel
  realization at the target rank (\cref{def:genuine-hinge-realization}),
  and if $G$ is simple, a generic one (\cref{def:rank-hypothesis}).
\end{lemma}
\begin{proof}
  \leanok
  \uses{lem:case-I-realization-nonsimple, lem:case-I-realization,
        lem:case-I-realization-h65, lem:generic-yields-genuine-hinge}
  If $G$ is not simple, \cref{lem:case-I-realization-nonsimple} (KT
  Lemma~6.2) applies directly. If $G$ is simple, case on whether some
  proper rigid subgraph $H$ with representative body $r$ has a simple
  contraction $G/E(H)$: if so, \cref{lem:case-I-realization} (KT
  Lemma~6.3) gives the generic realization; otherwise every contraction is
  non-simple and \cref{lem:case-I-realization-h65} (KT Lemma~6.5) gives
  it instead. Either way \cref{lem:generic-yields-genuine-hinge} then
  gives the nondegenerate-hinge form.
\end{proof}

\subsection{Case II: \texorpdfstring{$k > 0$}{k > 0} splitting (1-extension)}
\label{sec:molecular-algebraic-induction-caseII}

\begin{lemma}[Case II rank-lift accounting; the $+D$ core of KT Lemma 6.8]
  \label{lem:case-II-rank-lift}
  \lean{CombinatorialRigidity.Molecular.BodyHingeFramework.rankHypothesis_iff_finrank_pinnedMotions,
        CombinatorialRigidity.Molecular.PanelHingeFramework.rankHypothesis_withNormal_iff_finrank_pinnedMotions,
        CombinatorialRigidity.Molecular.BodyHingeFramework.pinnedMotions_le_withGraph,
        CombinatorialRigidity.Molecular.PanelHingeFramework.toBodyHinge_pinnedMotions_le_withGraph,
        CombinatorialRigidity.Molecular.BodyHingeFramework.pinnedMotions_withGraph_eq,
        CombinatorialRigidity.Molecular.PanelHingeFramework.toBodyHinge_pinnedMotions_withGraph_eq,
        CombinatorialRigidity.Molecular.BodyHingeFramework.hnew_of_isLink_incident,
        CombinatorialRigidity.Molecular.PanelHingeFramework.toBodyHinge_hnew_of_isLink_incident}
  \leanok
  \uses{def:rank-hypothesis, lem:rank-delete-vertex, def:framework-with-graph,
        lem:exists-independent-panel-extensor}
  In the codimension form of \cref{sec:molecular-rigidity-matrix},
  re-inserting (equivalently, pinning) a body $v$ shifts the realization
  count by exactly $D$. A framework $F$ realizes the target
  rank at $k'$ ($\dim Z(G,p) = D + k'$, \cref{def:rank-hypothesis})
  precisely when its body-$v$-pinned motion subspace has dimension $k'$.
  This is the algebraic $+D$ core of the panel-hinge $1$-extension: the
  pinned subspace is the null space of the rigidity matrix with $v$'s $D$
  columns deleted (the smaller graph $G - v$), and $\dim(\text{pinned at
  } v) + D = \dim Z(G,p)$ by the pin-a-body Lemma~5.1
  (\cref{lem:rank-delete-vertex}). Assembled on the panel layer: building
  the $1$-extension by choosing a panel normal $n$ for the re-inserted body
  $v$ (still unhinged, so the choice does not disturb the inductive null
  space, \cref{def:framework-with-graph}), the extended framework
  $P[v \mapsto n]^\flat$ realizes the target rank at $k'$ exactly when the
  splitting-off carries body-$v$-pinned motion dimension $k'$. The graph
  step of the $1$-extension is the unconditional inclusion: re-adding $v$'s
  two new hinge edges (passing from the splitting-off $G_v^{ab}$ up to $G$)
  can only \emph{shrink} the body-$v$-pinned motions, so the inductive
  realization of $G_v^{ab}$ bounds the extended framework's body-$v$-pinned
  dimension from above (the residual cut by the two new edges is the
  genericity-gated Claim~6.9 step). That inclusion is \emph{tight} exactly
  when every base-$v$-pinned motion of $G_v^{ab}$ already clears the two
  re-added $v$-incident hinge constraints $S a \in \operatorname{span}
  C(e_a)$, $S b \in \operatorname{span} C(e_b)$ (a base-pinned motion has
  $S v = 0$): under that hypothesis the body-$v$-pinned motions on $G$ and
  on $G_v^{ab}$ \emph{coincide}, so the extended framework's pinned
  dimension equals the inductive realization's and the $1$-extension lifts
  the rank by exactly $D$. The hypothesis is the honest content of
  Claim~6.9's general position --- supplied by placing the two new
  supporting extensors in general position
  (\cref{lem:exists-independent-panel-extensor}). Since the only links of
  $G$ outside $G_v^{ab}$ are the two $v$-incident edges $v a$, $v b$, and a
  base-pinned motion has $S v = 0$, each re-added hinge constraint
  $S v - S w \in \operatorname{span} C(e)$ collapses to the single
  span-membership $S w \in \operatorname{span} C(e)$ at the non-$v$
  endpoint, reducing the tightness hypothesis to exactly the two conditions
  $S a \in \operatorname{span} C(e_a)$, $S b \in \operatorname{span}
  C(e_b)$ that general position achieves.
\end{lemma}
\begin{proof}
  \leanok
  \uses{lem:rank-delete-vertex, def:framework-with-graph}
  Immediate from the pin-a-body identity
  $\dim(\text{pinned at } v) + D = \dim Z(G,p)$
  (\cref{lem:rank-delete-vertex}): substitute it into the definition of
  the rank hypothesis $\dim Z(G,p) = D + k'$ and cancel the $D$ term. For
  the panel assembly, the body-$v$-pinned motions of $P[v \mapsto n]^\flat$
  agree with those of $P^\flat$ when $v$ is unhinged
  (\cref{def:framework-with-graph}), so the $+D$ accounting applies to the
  pre-choice framework's pinned dimension.
\end{proof}

\begin{lemma}[Edge-substitution bridge for splitting-off]
  \label{lem:splitoff-edge-substitution}
  \lean{Graph.removeVertex_le, Graph.removeVertex_le_splitOff,
        Graph.isLink_incident_of_not_removeVertex}
  \leanok
  The splitting-off $G_v^{ab}$ is \emph{not} a subgraph of $G$: it deletes
  the two $v$-incident edges $v a$, $v b$ but adds a fresh short-circuit
  edge $e_0 \notin E(G)$ joining $a$ and $b$. The two graphs are instead an
  \emph{edge substitution} of one another, sharing the common subgraph
  $G - v$ (all of $G$ away from $v$): $G - v \le G$ and
  $G - v \le G_v^{ab}$. The second inclusion holds because every edge of
  $G - v$ is an edge of $G$, hence $\neq e_0$, so it survives into
  $G_v^{ab}$. Conversely, every link of $G$ \emph{lost} by passing to
  $G - v$ is incident to $v$: the re-added edges over the common lower
  bound are exactly the $v$-incident ones, which is what lets the genericity
  step (\cref{lem:case-II-rank-lift}) collapse each re-added hinge constraint
  on a $v$-pinned motion to a single span membership at its non-$v$ endpoint.
\end{lemma}
\begin{proof}
  \leanok
  Both inclusions are the definition of $\le$ on multigraphs (vertex-set
  and link inclusion). For $G - v \le G_v^{ab}$, a link of $G - v$ is a
  link $e\,x\,y$ of $G$ with $x, y \neq v$; since $e \in E(G)$ and
  $e_0 \notin E(G)$ we have $e \neq e_0$, so $e$ lands in $G_v^{ab}$'s
  $v$-avoiding branch.
\end{proof}

\begin{lemma}[Case II: splitting off a reducible vertex; KT Lemmas 6.7, 6.8]
  \label{lem:case-II}
  \lean{CombinatorialRigidity.Molecular.BodyHingeFramework.isInfinitesimallyRigidOn_insert_iff,
        CombinatorialRigidity.Molecular.PanelHingeFramework.rankHypothesis_withNormal_withGraph_iff_finrank_pinnedMotions}
  \leanok
  \uses{thm:minimal-kdof-reduction, lem:reduction-step,
        lem:rank-delete-vertex, lem:case-II-rank-lift, lem:cycle-realization,
        lem:splitoff-edge-substitution, lem:exists-independent-panel-extensor,
        lem:genericity-device}
  Let $G$ be a minimal $k$-dof-graph with $k > 0$ and a reducible
  degree-2 vertex $v$ (KT~4.6), and let $G_v^{ab}$ be the splitting-off
  at $v$ (a smaller minimal $k$-dof-graph by \cref{lem:reduction-step}).
  A realization of $G_v^{ab}$ at its inductive rank extends to a
  realization of $G$ at rank $D(|V(G)|-1) - k$. This is the panel-hinge
  analogue of Whiteley's bar-joint $1$-extension: re-inserting $v$ with
  its two hinges lifts the rank by exactly $D$ (Lemma~6.8, via the
  Claim~6.9 genericity argument).

  Concretely, with $P$ a panel framework on $G_v^{ab}$ in which $v$ is
  yet unhinged and $P.\mathrm{graph} \le G$, choosing a panel normal $n$
  for $v$ and enlarging the graph to $G$ gives the extended framework
  $(P[v \mapsto n]).\mathrm{withGraph}\,G$, which realizes the target
  rank at $k'$ \emph{iff} the splitting-off $P$ carries body-$v$-pinned
  dimension $k'$ --- the inductive realization lifts to $G$ with the two
  new hinge-row blocks accounting for the $+D$
  (\cref{lem:case-II-rank-lift}). The graph-side hypotheses are
  genericity-free: $v$ is unhinged in $G_v^{ab}$ (true before its two new
  edges are added) and every link of $G$ lost on passing to $G_v^{ab}$ is
  $v$-incident (\cref{lem:splitoff-edge-substitution}). The \emph{one}
  input from the genericity device (Phase~21b, Claim~6.9) is that each
  base-$v$-pinned motion lands in the two new edges' panel-support spans,
  $S a \in \operatorname{span} C(e_a)$, $S b \in \operatorname{span}
  C(e_b)$; this is false pointwise and holds only for the general-position
  normals the rank/dimension count selects
  (\cref{lem:exists-independent-panel-extensor}).
  As with Case~I, the $+D$ lift is accounted $\alpha$-independently through the
  bridge
  $\mathrm{isInfinitesimallyRigidOn\_insert\_iff}$: re-inserting the single
  body $v$ (so $V(G) = \{v\} \cup V(G_v^{ab})$), the framework is rigid on
  $V(G)$ \emph{iff} it is rigid on the splitting-off body set $V(G_v^{ab})$
  (the inductive realization) \emph{and} every motion pins $v$'s screw to the
  bodies it joins, $S v = S w$ for $w \in V(G_v^{ab})$. This references only
  $V(G)$, hence is $\alpha$-independent and composes through the induction. The
  $S v = S w$ condition is the honest geometric content of Claim~6.9's general
  position (the span-membership the device \cref{lem:genericity-device}
  supplies, via \cref{lem:exists-independent-panel-extensor}); the
  $1$-extension count itself ($v$'s two hinges add $D-1$ rows, $v$'s body adds
  one column block) is rank-side and genericity-free. The nullity-side iff is
  retained as the $\alpha$-dependent sibling.
\end{lemma}
\begin{proof}
  \leanok
  \uses{lem:rank-delete-vertex, lem:case-II-rank-lift,
        lem:splitoff-edge-substitution}
  Place $G_v^{ab}$ at its inductive realization; re-insert $v$ joined to
  $a$ and $b$ by hinges chosen generically. The rank-lift accounting
  \cref{lem:case-II-rank-lift} (via the pin-a-body Lemma~5.1,
  \cref{lem:rank-delete-vertex}) accounts for $v$'s screw freedoms through
  the unhinged-$v$ panel-normal invariance, and the genericity-gated
  tightness \cref{lem:case-II-rank-lift} (the $\ge$ half) makes the
  $1$-extension inclusion an equality once the span-membership input
  clears the two new $v$-incident edges (a base-pinned motion has
  $S v = 0$, so each constraint $S v - S w \in \operatorname{span} C(e)$
  collapses to $S w \in \operatorname{span} C(e)$ at the non-$v$
  endpoint). Hence the extended pinned dimension equals the inductive
  one and the rank lifts by exactly $D$, giving $D(|V|-1) - k$. The
  span-membership input is supplied by the genericity device (Claim~6.9,
  Phase~21b).
\end{proof}

\begin{lemma}[Case II realization, all $k \ge 0$: the splitting-off producer; KT Lemma 6.8]
  \label{lem:case-II-realization}
  \lean{CombinatorialRigidity.Molecular.PanelHingeFramework.case_II_realization_all_k}
  \leanok
  \uses{def:panel-hinge-framework, lem:case-II, lem:reduction-step-pos,
        lem:rank-polynomial-of-le-finrank}
  The generic-realization motive is closed under splitting off a
  degree-2 vertex at deficiency $k > 0$. By \cref{lem:reduction-step-pos}
  the splitting-off graph $G_v^{ab}$ is minimal $(k-1)$-dof; the all-$k$
  inductive hypothesis supplies a panel realization $Q$ with
  $\operatorname{rank} R(G_v^{ab}, Q) = D(|V|-2) - (k-1)$.  The
  deficiency-aware eq.-(6.12) placement brick places $v$'s normal at the
  shear $n_a + t\,n_b$ and produces a combined normal assignment $q_0$
  with $N = D(|V|-1) - k$ rows linearly independent in $R(G, q_0)$.  The
  W-suite shear argument (\cref{lem:case-III}) certifies one of the
  candidate $t$-values; the genericity device
  (\cref{lem:rank-polynomial-of-le-finrank}) re-realizes the attained rank
  at a general-position, algebraically-independent seed.
\end{lemma}
\begin{proof}
  \leanok
  \uses{lem:case-III, lem:rank-polynomial-of-le-finrank}
  The deficiency-aware placement brick (analogous to
  \cref{lem:case-II-realization-placement} at $k > 0$) supplies $N \le
  \operatorname{rank} R(G, q_0)$.  The W-suite (\cref{lem:case-III}) finds
  a good shear parameter; \cref{lem:rank-polynomial-of-le-finrank} lifts
  the rank lower bound to a nonzero rational polynomial whose non-roots
  realize deficiency $k$.
\end{proof}


\subsection{The genericity device (Claim 6.4/6.9)}
\label{sec:molecular-algebraic-induction-genericity}

The accounting nodes above (\cref{lem:case-I}, \cref{lem:case-II},
\cref{thm:theorem-55}, \cref{prop:rigidity-matrix-prop11}) and the
projective assembly of the cycle realization (\cref{lem:cycle-realization})
all reference one shared analytic crux, scoped out into the focused
sub-phase \textbf{Phase~21b} (the analytic sibling of the Phase-21a meet).
The device is the node \cref{lem:genericity-device}, now green.

\medskip\noindent\emph{What the device does and does not discharge.} The
device lifts a rank \emph{attained at one realization} to the same rank at a
generic point. It supplies the genericity inputs each accounting node carries
as a hypothesis (the block-triangular gluing inequality of \cref{lem:case-I},
the span membership of \cref{lem:case-II}, the generic-rank lower bound
$\mathtt{hgen}$ of \cref{prop:rigidity-matrix-prop11}, the projective assembly
of \cref{lem:cycle-realization}). It does \emph{not} build the realizations the
induction consumes. Two things stand between the device and a realization
producer: (i) the device must be \emph{applied} to a varying panel-coordinate
family, which needs the row coordinatization
\cref{lem:rows-polynomial-in-normals} (B0, now green) --- the prior phase only
ever invoked the device on a constant family; and (ii) the device needs a
\emph{seed} realization to lift. The device-to-motive closure itself --- lift a
seed's witnessed independent rows to a generic realization rigid on $V(G)$ --- is
the green glue \cref{lem:realization-of-independent-rows}, shared by both cases;
what remains in each is the \emph{placement}: the seed realization to feed it. For
Case~I the producer (\cref{lem:case-I-realization}) is now green-modulo: KT
eq.~(6.3)'s block-triangular rank-addition over one common framework supplies the
$D(|V(G)|-1)$ independent rows directly, its single residual being KT Claim~6.4
(\cref{lem:claim-6-4}, deferred to Phase~22b); for Case~II
the placement is the $1$-extension normal
(\cref{lem:case-II-realization-placement}), still red. Nor
does the device supply \cref{prop:rigidity-matrix-prop11}'s genericity-free
upper bound $\mathtt{hub}$, a distinct Phase-19-partition obligation carried as
an explicit hypothesis; the device underwrites only the $\mathtt{hgen}$ half.

\begin{lemma}[Genericity device, codimension form; KT Claim 6.4, Claim 6.9]
  \label{lem:genericity-device}
  \lean{CombinatorialRigidity.Molecular.exists_good_realization,
        CombinatorialRigidity.Molecular.exists_good_realization_const}
  \leanok
  \uses{def:rigidity-matrix, def:panel-hinge-framework}
  Fix a multigraph $G$ and regard a panel-hinge realization as a point
  $p \in \mathbb{R}^N$ of the panel-coordinate space (the per-vertex panel
  normals). Each entry of the panel-hinge rigidity matrix $R(G,p)$ is a
  polynomial in $p$ --- of degree two, bilinear in the normals
  (\cref{def:rigidity-matrix}, \cref{def:panel-hinge-framework}). Hence for
  each $r$ the locus $\{p : \operatorname{rank} R(G,p) \ge r\}$, cut out by
  the non-vanishing of some $r \times r$ minor, is Zariski-open, so
  $\operatorname{rank} R(G,p)$ is lower semicontinuous and attains its maximum
  on a dense generic set; a rank attained at one realization is attained
  generically. Dually, in the codimension convention of
  \cref{sec:molecular-rigidity-matrix} ($\operatorname{rank} R(G,p) = D|V| -
  \dim Z(G,p)$, \cref{def:dof-generic}), $\dim Z(G,p)$ is upper semicontinuous
  and minimized generically.
\end{lemma}
\begin{proof}
  \leanok
  A maximal non-vanishing minor of $R(G,p_0)$ is a polynomial in $p$, nonzero
  at $p_0$. Over the infinite field $\mathbb{R}$ a nonzero polynomial does not
  vanish identically (\texttt{MvPolynomial.funext}), so its non-vanishing
  locus is dense and nonempty; on that locus $\operatorname{rank} R(G,p)$ is at
  least the size of the minor. Several generic conditions intersect, each being
  the complement of a proper algebraic subset. See
  \cite[\S6.1, \S6.3]{katohTanigawa2011} (Claim~6.4 and its Case-II twin
  Claim~6.9).

  The formalization is genuinely \emph{multivariate}: the parameter ranges
  over the full panel-coordinate space $\mathbb{R}^N$, not a single affine line.
  The minor's Gram determinant is one multivariate polynomial; nonzero at the
  witness realization, it has a non-root by \texttt{MvPolynomial.funext}, and
  there the rank is at least the witnessed corank. In the codimension shape the
  consumers carry, this is $\#s + \dim Z(F\,p) \le D|V|$ at a produced
  realization. The constant-family closure
  (\texttt{exists\_good\_realization\_const}) reads the bound off a single
  hand-built realization, the form Case~I's block-triangular gluing consumes.
\end{proof}

\noindent The device thus discharges the genericity \emph{hypotheses} the
accounting nodes carry, but not the realization-layer constructions that feed
the induction. With the row coordinatization \cref{lem:rows-polynomial-in-normals}
green and the Case~I producer (\cref{lem:case-I-realization}) now green-modulo KT
Claim~6.4 (\cref{lem:claim-6-4}, deferred to Phase~22b), the still-red
constructions are that residual Claim~6.4 and the Case~II $1$-extension
realization (\cref{lem:case-II-realization}). The device is also consumed by the
Case~III candidate-framework genericity of Phases~22--23.

\subsection{The \texorpdfstring{$V(G)$}{V(G)}-relative count (device output to relative motive)}
\label{sec:molecular-algebraic-induction-relative}

The device (\cref{lem:genericity-device}) produces a generic realization at
which the rigidity-matrix corank is minimized, but its bound is stated
\emph{absolutely}, over the whole screw-assignment space $\R^{D|V_\alpha|}$ of
the ambient body type: a produced point satisfies
$\#s + \dim Z(G,p) \le D|V_\alpha|$. The realization motive
(\cref{def:rank-hypothesis}), by contrast, is $V(G)$-relative
($\mathrm{IsInfinitesimallyRigidOn}\,V(G)$). For a non-spanning inductive
subgraph the two differ by the $D(|V_\alpha| - |V(G)|)$ free dimensions of the
bodies outside $V(G)$ --- each carries no hinge constraint, hence is a free
motion --- so the absolute count alone never reaches the relative motive. The
three nodes here are the adapter the producers (\cref{lem:case-I-realization},
\cref{lem:case-II-realization}) compose with the device to land the relative
motive; they are the ``rank-equality bridge'' \cref{def:rank-hypothesis}
anticipates.

\begin{lemma}[Isolated bodies carry no rigidity: the relative split]
  \label{lem:relative-screw-split}
  \lean{CombinatorialRigidity.Molecular.BodyHingeFramework.finrank_pinnedMotionsOn_vertexSet}
  \leanok
  \uses{def:rigidity-matrix, def:rank-hypothesis, def:pinned-motions-on}
  The screw-assignment space $\R^{D|V_\alpha|} = (\alpha \to \mathrm{Screw})$
  splits as a product along $V(G)$ and its complement, and a body
  $w \in \alpha \setminus V(G)$ carries no hinge constraint, so its $D$ screw
  coordinates are free in every infinitesimal motion. The block pin on the whole
  vertex set $V(G)$ therefore captures exactly the free isolated bodies:
  \[
    \dim\bigl(\mathrm{pinnedMotionsOn}\,V(G)\bigr) \;=\; D\,(|V_\alpha| - |V(G)|),
  \]
  the $\alpha$-dependent free summand of $\dim Z(G,p)$. Combined with the
  block-pin lower bound \cref{lem:pinned-motions-on-rank-bound}, this is what
  isolates the $\alpha$-independent residual rank.
\end{lemma}
\begin{proof}
  \leanok
  An assignment vanishing on $V(G)$ automatically satisfies every hinge
  constraint $S u - S v = 0$ (both endpoints lie in $V(G)$), so
  $\mathrm{pinnedMotionsOn}\,V(G)$ is the kernel of the projection onto the
  $V(G)$ coordinates, $\bigsqcap_{i \in V(G)} \ker(\mathrm{proj}_i)$. Mathlib's
  \texttt{iInfKerProjEquiv} identifies that kernel with the product over the
  complement $\alpha \setminus V(G)$, of dimension $D(|V_\alpha|-|V(G)|)$ by
  \texttt{finrank\_pi\_const}.
\end{proof}

\begin{lemma}[The genericity device, $V(G)$-relative form]
  \label{lem:relative-device-count}
  \lean{CombinatorialRigidity.Molecular.PanelHingeFramework.exists_relative_full_count_ofParam}
  \leanok
  \uses{lem:genericity-device, lem:relative-screw-split,
        lem:rows-polynomial-in-normals}
  Substituting the relative split \cref{lem:relative-screw-split} for the
  absolute screw-count in the device \cref{lem:genericity-device} (applied to
  the varying panel family via the row coordinatization
  \cref{lem:rows-polynomial-in-normals}, B0), a generic panel assignment with at
  least $\#s \ge D(|V(G)|-1)$ witnessed independent rows attains
  \[
    \dim Z(G,p) \;\le\; D\,\bigl(|\alpha \setminus V(G)| + 1\bigr),
  \]
  the device's absolute codimension bound with the $D(|V(G)|-1)$ rows cancelled
  --- i.e.\ the relative full count, the ambient $D|\alpha \setminus V(G)|$ free
  dimensions plus the $D$ trivial-motion dimensions and no residual relative
  corank. Mechanical given \cref{lem:relative-screw-split}: re-wrap the device's
  conclusion using $|\alpha| = |V(G)| + |\alpha \setminus V(G)|$.
\end{lemma}
\begin{proof}
  \leanok
  The device (\cref{lem:rows-polynomial-in-normals}) gives
  $\#s + \dim Z(G,p) \le D|\alpha|$ at the produced point; substitute
  $|\alpha| = |V(G)| + |\alpha \setminus V(G)|$ and the count
  $\#s \ge D(|V(G)|-1)$, then cancel.
\end{proof}

\begin{lemma}[Relative count $\Rightarrow$ the relative motive]
  \label{lem:isInfRigidOn-of-relative-count}
  \lean{CombinatorialRigidity.Molecular.BodyHingeFramework.isInfinitesimallyRigidOn_vertexSet_of_finrank_le}
  \leanok
  \uses{lem:relative-device-count, lem:relative-screw-split,
        lem:pinned-motions-on-rank-bound, lem:rank-delete-vertex, lem:case-I}
  A realization at the relative full count ($k'=0$ in
  \cref{lem:relative-device-count}, i.e.\ $\dim Z(G,p) \le
  D(|\alpha \setminus V(G)| + 1)$) is infinitesimally rigid on $V(G)$:
  $\mathrm{IsInfinitesimallyRigidOn}\,V(G)$. This is the adapter the producers
  compose with the device output, routed through the relative split
  \cref{lem:relative-screw-split} and the relativized Case~I bridge
  \cref{lem:case-I}
  ($\mathrm{isInfinitesimallyRigidOn\_iff\_pinnedMotionsOn\_le}$).
\end{lemma}
\begin{proof}
  \leanok
  Pick any $v_0 \in V(G)$ and read rigidity-on-$V(G)$ off the Case~I bridge
  \cref{lem:case-I} at the trivially-rigid singleton block $\{v_0\}$: it reduces
  to $\mathrm{pinnedMotionsOn}\,\{v_0\} \le \mathrm{pinnedMotionsOn}\,V(G)$. The
  reverse containment is automatic, so it suffices to match dimensions. The
  single-body pin has $\dim(\mathrm{pinnedMotions}\,v_0) = \dim Z(G,p) - D$
  (\cref{lem:rank-delete-vertex}), capped at $D|\alpha \setminus V(G)|$ by the
  hypothesis; the block pin on $V(G)$ has exactly that dimension
  (\cref{lem:relative-screw-split}). Equal dimensions on a containment force the
  needed equality.
\end{proof}

\begin{lemma}[Realization glue from a witnessed independent row family]
  \label{lem:realization-of-independent-rows}
  \lean{CombinatorialRigidity.Molecular.PanelHingeFramework.hasFullRankRealization_of_independent_panelRow}
  \leanok
  \uses{def:rank-hypothesis, lem:relative-device-count,
        lem:isInfRigidOn-of-relative-count, lem:rows-polynomial-in-normals}
  The honest device closure both case producers reduce to once their geometry is
  placed. Fix a multigraph $G$ on a nonempty body set $V(G)$ and an endpoint
  selector $\mathrm{ends}$ recording each edge's link. If at \emph{one} free
  normal assignment $q_0$ the rigidity rows of the free-normal panel family
  $\mathrm{ofNormals}\,G\,\mathrm{ends}\,q_0$ indexed by $s$ are linearly
  independent and $\#s \ge D(|V(G)|-1)$, then $G$ has a full-rank panel
  realization ($\mathrm{HasFullRankRealization}$, the $V(G)$-relative rank).
  This is the composition $\text{N2} \circ \text{N3}$: the device's relative
  count \cref{lem:relative-device-count} lifts the witnessed corank at $q_0$ to a
  generic normal assignment $q$, and the relative-count adapter
  \cref{lem:isInfRigidOn-of-relative-count} turns that into rigidity on
  $V(G) = V(\mathrm{ofNormals}\,G\,\mathrm{ends}\,q)$. The witnessed independent
  family $(q_0, s)$ is the \emph{satisfiable} geometric input the placement
  supplies, not the rank the lemma concludes --- the honesty split matching the
  Case~I splice glue \cref{lem:case-I-splice-seed}.
\end{lemma}
\begin{proof}
  \leanok
  Apply \cref{lem:relative-device-count} to the witness $(q_0, s)$ to obtain a
  normal assignment $q$ with $\dim Z(G,q) \le D(|\alpha \setminus V(G)| + 1)$,
  then \cref{lem:isInfRigidOn-of-relative-count} on
  $\mathrm{ofNormals}\,G\,\mathrm{ends}\,q$ (whose vertex set is $V(G)$).
\end{proof}

\begin{lemma}[Case II placement: the $1$-extension normal; KT (6.12)]
  \label{lem:case-II-realization-placement}
  \uses{def:panel-hinge-framework, lem:case-II, lem:genericity-device,
        lem:exists-independent-panel-extensor,
        lem:case-II-placement-new-rows, lem:case-II-placement-old-rows,
        lem:case-II-placement-old-rows-extract,
        lem:case-II-placement-block-independent,
        lem:case-II-placement-e0-recovery}
  The remaining geometric obligation of \cref{lem:case-II-realization}: from the
  inductive realization of the splitting-off $G_v^{ab}$, construct a free normal
  assignment $q_0$ for the parent graph $G$ --- in particular a panel normal $n$
  for the re-inserted body $v$ --- together with a linearly independent family of
  $D(|V(G)|-1)$ rigidity rows of $\mathrm{ofNormals}\,G\,\mathrm{ends}\,q_0$.
  Geometrically: re-insert $v$ joined to its two neighbours $a, b$ by hinges
  whose supporting extensors are in general position
  (\cref{lem:exists-independent-panel-extensor}), so the two new edges contribute
  $+(D-1)$ independent rows that, combined with the inductive realization's rows
  threaded through the common subgraph $G - v$ (recovering the short-circuit edge
  $e_0$'s constraint from $v$'s two new edges --- the edge substitution
  $G_v^{ab}$ adds $e_0$ but deletes $v$'s edges, so the count is \emph{not} a free
  transport), witness the full target count. This is the shallower parallel of
  the Case~I splice placement (one re-inserted body, $+(D-1)$ rows). \emph{Red}
  (Phase~21b realization layer): the witnessed-independent-row construction across
  the edge substitution is the genuine content; once green it feeds the glue
  \cref{lem:realization-of-independent-rows}.

  \medskip\noindent\emph{Decomposition (Phase-21b realization-layer plan, revised
  2026-06-04).} Per the hand-off discipline for a genuinely-geometric node, this
  placement breaks into buildable sub-nodes plus a closing assembly, mirroring the
  Case~I splice placement's boundary-panel/block-triangular split. The seed family
  $s$ partitions into the \emph{new-edge block} (the two hinges re-attaching $v$)
  and the \emph{old block} (the inductive realization's rows, read off the common
  subgraph $G - v$); the count is $(D-1) + D(|V(G)|-2) = D(|V(G)|-1)$.
  \begin{enumerate}
    \item \cref{lem:case-II-placement-new-rows} (N7b-1, green): choose the normal $n$
      for $v$ so that $v$'s two new hinges supply $D-1$ linearly independent panel rows.
    \item \cref{lem:case-II-placement-old-rows-extract} (N7b-0, green): \emph{produce}
      the old block --- extract $D(|V(G)|-2)$ independent panel rows from the inductive
      realization's rigidity on $V(G_v^{ab})$ (the rank-equals-codimension forward
      direction, the converse of \cref{lem:isInfRigidOn-of-relative-count}).
    \item \cref{lem:case-II-placement-old-rows} (N7b-2, green): transport the
      $e_0$-\emph{free} subfamily of that block onto $G$ through the common subgraph
      $G - v$.
    \item \cref{lem:case-II-placement-block-independent} (N7b-3, green): the two blocks
      are jointly independent (the new-edge rows live in the $v$-column block, the
      old rows in its complement; pin-a-body Lemma~5.1,
      \cref{lem:rank-delete-vertex}, makes the union block-triangular).
    \item \cref{lem:case-II-placement-e0-recovery} (N7b-4, \emph{superseded}): a
      scrutiny pass found N7b-0/1/2/3 do not compose --- N7b-0 extracts an old block
      with no control over $e_0$-rows, but N7b-2 transports only the $e_0$-free rows,
      so the count can fail --- and a recon pass then found the row-side fix
      (``produce an $e_0$-free old block of full count'') \emph{geometrically
      unbuildable} ($G - v$ is not rigid, so no such block exists in $G_v^{ab}$). The
      genuine rank-lift is re-framed \emph{motion-side} (M1--M3,
      \cref{lem:case-II-placement-motion-side-assembly}), which supersedes this whole
      row-side assembly path for Case~II. See ``The rank-lift, motion-side'' below.
  \end{enumerate}
  \emph{Superseded (2026-06-04):} the row-side assembly below was the original plan,
  but the missing producer N7b-4 is geometrically unbuildable in row-side form; the
  Case~II producer instead routes through the motion-side re-frame
  (\cref{lem:case-II-placement-motion-side-assembly}, M3), which concludes rigidity
  on $V(G)$ directly without a witnessed independent row family. The row-side
  sub-nodes N7b-0/1/2/3 + glue N7a stay green and are reused by Case~I.
\end{lemma}
\begin{proof}
  \uses{lem:case-II, lem:exists-independent-panel-extensor,
        lem:case-II-placement-new-rows, lem:case-II-placement-old-rows,
        lem:case-II-placement-old-rows-extract,
        lem:case-II-placement-block-independent,
        lem:case-II-placement-e0-recovery}
  Take $q_0$ to agree with the inductive realization's normals on
  $V(G)\setminus\{v\} = V(G_v^{ab})$ and to assign $v$ the general-position normal
  $n$ of \cref{lem:case-II-placement-new-rows}. Let $s$ be the union of the
  new-edge block (\cref{lem:case-II-placement-new-rows}, $D-1$ rows) and the
  $e_0$-free old block (\cref{lem:case-II-placement-e0-recovery,
  lem:case-II-placement-old-rows}, $D(|V(G)|-2)$ rows);
  \cref{lem:case-II-placement-block-independent} makes the union independent, and
  the count is $D(|V(G)|-1)$. The genuine content is the edge substitution: $s$ is
  \emph{not} a free transport of $G_v^{ab}$'s rows (the short-circuit $e_0$ is
  deleted in $G$, $v$'s two edges added), so the old block is recovered through the
  common subgraph $G - v$ ($\mathrm{removeVertex}$), not along an inclusion
  $G_v^{ab}\le G$; the $e_0$-row's constraint comes from $v$'s two new edges
  (\cref{lem:case-II-placement-e0-recovery}). Deferred to the Phase-21b realization
  layer; see \cite[\S6.3]{katohTanigawa2011} (Lemma~6.8).
\end{proof}

\begin{lemma}[N7b-1: a re-inserted body's transversal hinge gives $D-1$ independent panel rows]
  \label{lem:case-II-placement-new-rows}
  \lean{CombinatorialRigidity.Molecular.BodyHingeFramework.exists_independent_panelRow_of_edge,
        CombinatorialRigidity.Molecular.BodyHingeFramework.exists_independent_panelRow_subfamily_of_edge}
  \leanok
  \uses{def:panel-hinge-framework, def:hinge-row-block, def:rigidity-matrix, lem:rank-delete-vertex}
  Re-inserting the degree-2 body $v$ joins it to its neighbours $a, b$ by two hinges. Each such
  hinge $e = uv$ that is \emph{transversal} --- distinct endpoints and a nonzero supporting extensor
  $C(p(e))$, the general-position condition supplied by choosing $v$'s normal $n$ off the
  support-extensor locus (\cref{lem:exists-independent-panel-extensor}) --- contributes a linearly
  independent family of $D-1 = \mathrm{screwDim}\,k - 1$ rigidity rows of
  $\mathrm{ofNormals}\,G\,\mathrm{ends}\,q_0$, each a member of \emph{that single edge's} panel-row
  span $\langle \mathrm{panelRow}\,\mathrm{ends}\,(e, \cdot, \cdot)\rangle$. These are the $+(D-1)$
  rows the $1$-extension adds in $v$'s column screw block: the hinge-row block $r(p(e))$ is
  $(D-1)$-dimensional (\cref{lem:rank-delete-vertex}, the rank-count input), and its basis lifts
  through the relative-screw evaluation $\mathrm{hingeRow}\,u\,v$ into the per-edge panel-row span.
  This is the panel-row form of the per-edge brick of \cref{def:rigidity-matrix}
  ($\mathrm{exists\_independent\_rigidityRows\_of\_edge}$), restricted to
  membership in this edge's panel rows so the placement assembly
  (\cref{lem:case-II-realization-placement}) can thread it into the device-consuming
  $\mathrm{panelRow}$ family of \cref{lem:realization-of-independent-rows}.

  \medskip\noindent\emph{Honest index-subfamily form
  ($\mathrm{exists\_independent\_panelRow\_subfamily\_of\_edge}$).} The glue
  \cref{lem:realization-of-independent-rows} consumes not a family of \emph{members} of the per-edge
  span but a literal $\mathrm{panelRow}\,\mathrm{ends}$-subfamily indexed by a \emph{set} of panel-row
  indices. This companion form bridges that gap: a transversal hinge $e$ yields an index subset
  $s \subseteq \{e\}\times(\Lambda^k\text{-pairs})$ of size $D-1$ whose \emph{actual}
  $\mathrm{panelRow}\,\mathrm{ends}$-subfamily is independent --- the honest input
  \cref{lem:case-II-realization-placement} threads into the genericity device.
\end{lemma}
\begin{proof}
  \leanok
  \uses{def:hinge-row-block, lem:rank-delete-vertex}
  The transversal hinge's row block $r(p(e)) = (\mathrm{span}\,C(p(e)))^\perp$ has dimension $D-1$
  (\cref{lem:rank-delete-vertex}); pick a basis of it. Its image under the relative-screw evaluation
  $\mathrm{hingeRow}\,u\,v = (\mathrm{screwDiff}\,u\,v)^\ast$ stays independent (the dual map is
  injective at distinct endpoints) and, by the per-edge span identity
  $\mathrm{span\_panelRow\_edge\_eq}$ (the annihilator-family span
  $\mathrm{span\_annihRow\_eq\_dualAnnihilator}$ carried through $\mathrm{hingeRow}\,u\,v$), each row
  lies in the single edge's panel-row span. For the index-subfamily form, the per-edge family
  $(t_1,t_2)\mapsto\mathrm{panelRow}\,\mathrm{ends}\,(e,t_1,t_2)$ spans a $(D-1)$-dimensional space, so
  extracting an independent subfamily of its generators at the span's $\mathrm{finrank}$
  ($\mathrm{Submodule.exists\_fun\_fin\_finrank\_span\_eq}$) and re-indexing each chosen row by its
  $\Lambda^k$-pair gives the honest index subset. See \cite[\S6.3]{katohTanigawa2011}.
\end{proof}

\begin{lemma}[N7b-2: the inductive rows transport through the common subgraph $G-v$]
  \label{lem:case-II-placement-old-rows}
  \lean{CombinatorialRigidity.Molecular.PanelHingeFramework.exists_independent_panelRow_transport}
  \leanok
  \uses{def:panel-hinge-framework, lem:case-II, lem:rows-polynomial-in-normals,
        lem:case-II-placement-old-rows-extract}
  Given the inductive realization's $D(|V(G)|-2)$ independent rows of
  $\mathrm{ofNormals}\,G_v^{ab}\,\mathrm{ends}'\,q_0'$ (produced by
  \cref{lem:case-II-placement-old-rows-extract} from the inductive rigidity),
  transport them to $D(|V(G)|-2)$ independent rows of
  $\mathrm{ofNormals}\,G\,\mathrm{ends}\,q_0$ on $G$. Concretely: along an \emph{injective} reindex $f : s_2 \to s_1$ selecting the
  $e_0$-free subfamily, with each selected row matching across the graph swap (the per-row
  hypothesis $\mathrm{hrow}$), the family is again linearly independent as rows of
  $\mathrm{ofNormals}\,G\,\mathrm{ends}\,q_0$. The transport is \emph{not} along an
  inclusion (neither $G_v^{ab}$ nor $G$ is a subgraph of the other): both sit above the
  common subgraph $G - v$ ($\mathrm{removeVertex}\,v$; $\mathrm{removeVertex\_le}$ and
  $\mathrm{removeVertex\_le\_splitOff}$, green), and the inductive rows that avoid $v$
  --- equivalently avoid the short-circuit edge $e_0$ --- are exactly the rows of
  $G - v$, which survive into $G$ unchanged. The per-row match is where the common
  subgraph enters: a panel row of $\mathrm{ofNormals}$ reads only $\mathrm{ends}$ and the
  normal assignment, not the underlying graph
  (\cref{lem:rows-polynomial-in-normals}, via $\mathrm{toBodyHinge\_supportExtensor}$), so
  when the assembly (\cref{lem:case-II-realization-placement}) picks $q_0$ extending the
  inductive normals and $\mathrm{ends}$ agreeing on $G - v$'s edges, the match is
  definitional. The short-circuit edge $e_0$'s constraint is \emph{not} transported here;
  it is recovered from $v$'s two new edges in N7b-1.
\end{lemma}
\begin{proof}
  \leanok
  \uses{lem:case-II, lem:rows-polynomial-in-normals}
  Restrict the inductive independent family to the rows indexed by edges of
  $G - v$ (drop the $e_0$-row), and read them as rows of $G$ through
  $\mathrm{removeVertex\_le}$. Independence is inherited as a subfamily of an
  independent family ($\mathrm{LinearIndependent.comp}$ along the injective reindex,
  the family equality given by the per-row match). See
  \cite[\S6.3]{katohTanigawa2011} (Lemma~6.8).
\end{proof}

\begin{lemma}[N7b-0: a rigid realization carries a full-rank independent row family]
  \label{lem:case-II-placement-old-rows-extract}
  \lean{CombinatorialRigidity.Molecular.BodyHingeFramework.exists_independent_panelRow_subfamily_of_rigidOn,
        CombinatorialRigidity.Molecular.BodyHingeFramework.finrank_infinitesimalMotions_of_isInfinitesimallyRigidOn_vertexSet}
  \leanok
  \uses{def:panel-hinge-framework, def:rigidity-matrix, def:rank-hypothesis,
        lem:rank-delete-vertex, lem:relative-screw-split}
  The \emph{producer} of the old block that the transport
  \cref{lem:case-II-placement-old-rows} consumes: from the inductive realization
  of $G_v^{ab}$ --- a panel-hinge framework infinitesimally rigid on its own
  vertex set $V(G_v^{ab}) = V(G)\setminus\{v\}$ ($\mathrm{HasFullRankRealization}$,
  the realization motive) --- extract a linearly independent family of
  $D(|V(G_v^{ab})|-1) = D(|V(G)|-2)$ \emph{actual} $\mathrm{panelRow}$ rows of
  that realization (in the honest index-subfamily form
  \cref{lem:case-II-placement-new-rows} supplies for the new block). This is the
  \emph{forward converse} of the relative-count adapter
  \cref{lem:isInfRigidOn-of-relative-count}: that node turns a witnessed corank
  $\#s \ge D(|V|-1)$ \emph{into} rigidity on $V(G)$; this node runs the
  implication backward --- rigidity on $V(G_v^{ab})$ forces
  $\dim Z = D(|V_\alpha| - |V(G_v^{ab})| + 1)$
  (\cref{lem:relative-screw-split}), so $\mathrm{rank}\,R = D|V_\alpha| - \dim Z =
  D(|V(G_v^{ab})|-1)$, and the rigidity rows, spanning a space of that dimension,
  contain an independent subfamily of size $D(|V(G_v^{ab})|-1)$
  ($\mathrm{Submodule.exists\_fun\_fin\_finrank\_span\_eq}$ on the rigidity-row
  span, re-indexed into a $\mathrm{panelRow}$-index subset as in
  \cref{lem:case-II-placement-new-rows}'s honest form). This
  rank-equals-codimension forward direction is the genuine deficit between
  ``rigid on a set'' and ``carries that many independent rows'' --- it is
  \emph{not} discharged by the transport \cref{lem:case-II-placement-old-rows}
  (which only carries an \emph{already-witnessed} independent family across the
  graph swap) nor by the graph-forest existence \cref{def:rigidity-matrix}
  ($\mathrm{exists\_independent\_rigidityRows\_of\_forest}$, which sources rows
  from a transversal-hinge spanning forest, of count bounded by the forest size
  rather than the full $D(|V|-1)$ unless the block is itself rigid).
\end{lemma}
\begin{proof}
  \leanok
  \uses{lem:rank-delete-vertex, lem:relative-screw-split}
  The realization is rigid on $V(G_v^{ab})$, so $\mathrm{pinnedMotionsOn}\,
  V(G_v^{ab})$ pins all residual freedom: its dimension is
  $D(|V_\alpha| - |V(G_v^{ab})|)$ (\cref{lem:relative-screw-split}), and the
  single-body pin gives $\dim Z = D + D(|V_\alpha| - |V(G_v^{ab})|)$
  (\cref{lem:rank-delete-vertex}), whence $\mathrm{rank}\,R = D|V_\alpha| - \dim Z
  = D(|V(G_v^{ab})|-1)$. The rigidity-row span has exactly that dimension, so
  $\mathrm{Submodule.exists\_fun\_fin\_finrank\_span\_eq}$ extracts an independent
  subfamily of that many actual rows; the $\mathrm{panelRow}$-index re-indexing is
  as in \cref{lem:case-II-placement-new-rows}'s honest form. See
  \cite[\S6.3]{katohTanigawa2011} (Lemma~6.8).
\end{proof}

\begin{lemma}[N7b-3: the new-edge and old blocks are jointly independent (pin-a-body column split)]
  \label{lem:case-II-placement-block-independent}
  \lean{CombinatorialRigidity.Molecular.BodyHingeFramework.linearIndependent_sum_pinned_block}
  \leanok
  \uses{def:panel-hinge-framework, lem:rank-delete-vertex}
  The pin-a-body column decomposition (Lemma~5.1, \cref{lem:rank-delete-vertex})
  certifying that the new-edge block (\cref{lem:case-II-placement-new-rows}, the
  rows carried by the re-inserted body $v$'s incident hinges, in the $v$-column
  screw block) and the old block (\cref{lem:case-II-placement-old-rows}, the rows
  of the common subgraph $G-v$, supported off $v$'s columns) form a jointly
  linearly independent family. The split is the body-$v$ screw column: a screw
  assignment pinned to $v$ alone ($\mathrm{Function.update}\,0\,v\,x$) reads the
  old rows as $0$ (their edges avoid $v$) while reading the new rows through $v$'s
  screw column $\mathrm{rn}\,i \circ \mathrm{single}_v$. So if the new rows are
  independent \emph{as functionals of $v$'s screw} (the $D-1$ column-block rows of
  N7b-1) and the old rows are independent (the $D(|V(G)|-2)$ inductive rows of
  N7b-2), their union $\mathrm{Sum.elim}\,\mathrm{rn}\,\mathrm{ro}$ is independent.
  This is the panel-row form of the $1$-extension's exact $+D$ rank lift the
  Case~II accounting \cref{lem:case-II} pins; the assembly
  \cref{lem:case-II-realization-placement} (N7b) supplies the two blocks (whence
  the count $D(|V(G)|-1)$) and discharges the column-split hypotheses by
  construction.
\end{lemma}
\begin{proof}
  \leanok
  \uses{lem:rank-delete-vertex}
  A vanishing combination of the union, evaluated at the body-$v$ pin
  $\mathrm{Function.update}\,0\,v\,x$, collapses to the new block alone (the old
  terms vanish, edges $e_0$-free away from $v$), forcing the new coefficients to
  vanish by their pinned independence; the residual is then a vanishing
  combination of the old block, forcing the old coefficients by their
  independence (pin-a-body Lemma~5.1, \cref{lem:rank-delete-vertex}). See
  \cite[\S6.3]{katohTanigawa2011}.
\end{proof}

\subsubsection*{The rank-lift, motion-side (the N7b-4 re-frame, 2026-06-04)}

A scrutiny+recon pass on the N7b assembly (Phase 21b, 2026-06-04) found the
\emph{row-side} framing of the rank-lift below
(\cref{lem:case-II-placement-e0-recovery}, ``produce an $e_0$-free old block of
full count'') is \emph{geometrically wrong}, and re-frames the genuine content
\emph{motion-side}. The finding, the corrected statement, and the buildable
sub-node decomposition follow; the row-side node is kept (struck) for the
audit trail.

\medskip\noindent\emph{Why the row-side framing fails.} An $e_0$-free old block
of $G_v^{ab}$ is a family of rigidity rows indexed by edges of $G_v^{ab}$ other
than $e_0$ --- i.e.\ rows of the common subgraph $G - v$. But $G - v$ is
\emph{not} rigid on $V(G)\setminus\{v\}$: it is missing both the short-circuit
$e_0$ and $v$'s two edges, so its rows do \emph{not} span the full
$D(|V(G)|-2)$-dimensional rigidity-row space of the rigid $G_v^{ab}$. Equivalently,
the $e_0$-row of $G_v^{ab}$ is genuinely needed in $G_v^{ab}$ --- it is \emph{not}
redundant there --- so no $e_0$-free old block of the full count exists in
$G_v^{ab}$ alone. The $e_0$-row only becomes recoverable \emph{in $G$}, after $v$'s
two new hinge rows are present; the row-counting bookkeeping that the green
sub-nodes N7b-0/1/2/3 set up cannot express this, because it asks for the old
block before $v$ is re-attached.

\medskip\noindent\emph{The corrected, motion-side rank-lift.} The realization
motive is $V(G)$-relative ($\mathrm{IsInfinitesimallyRigidOn}\,V(G)$,
\cref{def:rank-hypothesis}), so the producer is most directly a
\emph{motion-space} statement, in the established idiom of the Case~I splice
\emph{glue} \cref{lem:case-I-splice-seed} (re-adding edges only shrinks the
motion space, \cref{lem:motions-mono-of-graph-le}): \emph{if $G_v^{ab}$ is
rigid on $V(G)\setminus\{v\}$ and $v$'s two re-attaching hinges $e_a = va$,
$e_b = vb$ are transversal with $C(e_a), C(e_b)$ linearly independent, then $G$ is
rigid on $V(G)$.} The argument bypasses the $e_0$ bookkeeping entirely:
\begin{enumerate}
  \item It suffices to be rigid on $V(G)$ for the subgraph
    $G'' := (G-v) \cup \{e_a, e_b\} \le G$ (re-adding $G$'s remaining edges only
    shrinks the motion space, \cref{lem:motions-mono-of-graph-le}); $G''$ shares
    $G - v$ with $G_v^{ab}$.
  \item Let $S$ be an infinitesimal motion of $G''$ constant on
    $V(G)\setminus\{v\}$ (the hypothesis of rigid-on); in particular $S a = S b$.
    On $G - v$ this is also a motion of $G_v^{ab}$, so rigidity of $G_v^{ab}$ gives
    $S$ constant on $V(G)\setminus\{v\}$ (already assumed) --- the inductive input.
  \item The two new hinge constraints give
    $S v - S a \in \operatorname{span} C(e_a)$ and
    $S v - S b = S v - S a \in \operatorname{span} C(e_b)$, so
    $S v - S a \in \operatorname{span} C(e_a) \cap \operatorname{span} C(e_b)$.
    With $C(e_a), C(e_b)$ \emph{linearly independent}, the two $1$-dimensional spans
    meet only at $0$, so $S v = S a$: $v$ is pinned, $S$ is constant on $V(G)$.
\end{enumerate}
The general-position input ``$C(e_a), C(e_b)$ linearly independent'' is exactly the
$m = 2 \le D$ case of \cref{lem:exists-independent-panel-extensor}
($\mathrm{exists\_independent\_panelSupportExtensor}$), realized by choosing $v$'s
panel normal off the relevant locus --- the same satisfiable general-position
output N7b-1 (\cref{lem:case-II-placement-new-rows}) already consumes for the new
block. This is the genuine geometric heart of why a $1$-extension preserves rank
(KT~(6.12)), and unlike the row-side framing it is \emph{provable}: the meet step
is the elementary ``two distinct lines through the origin intersect trivially'',
not a deep extensor computation.

\medskip\noindent\emph{Buildable sub-node decomposition (N7b-4, motion-side).}
\begin{enumerate}
  \item[(M1)] \cref{lem:case-II-placement-disjoint-line-meet}: $C, C' \ne 0$ linearly
    independent $\Rightarrow \operatorname{span}\{C\} \cap \operatorname{span}\{C'\}
    = 0$ (general fact about $1$-dimensional spans; the meet step). Pure linear
    algebra, mathlib-supported.
  \item[(M2)] \cref{lem:case-II-placement-pin-vertex}: a $G''$-motion constant on
    $V(G)\setminus\{v\}$, with $v$ joined to $a, b$ by transversal hinges with
    independent support extensors, is constant on all of $V(G)$ (the pin, via M1 +
    the two hinge constraints). The crux.
  \item[(M3)] \cref{lem:case-II-placement-motion-side-assembly}: assemble M2 with the
    subgraph monotonicity \cref{lem:motions-mono-of-graph-le} and the inductive
    realization's rigidity on $V(G)\setminus\{v\}$ to conclude $G$ rigid on $V(G)$;
    feed the resulting $\mathrm{HasFullRankRealization}\,G$ to
    \cref{lem:case-II-realization}. The general-position normal for $v$ comes from
    \cref{lem:exists-independent-panel-extensor} ($m = 2$).
\end{enumerate}
This re-frame supersedes both the N7b row-side assembly path
(\cref{lem:case-II-realization-placement}) and the device-closure glue
\cref{lem:realization-of-independent-rows} \emph{for Case~II}: the motion-side
producer concludes $\mathrm{IsInfinitesimallyRigidOn}\,V(G)$ directly from the
inductive rigidity, with no detour through a witnessed independent row family.
(The green row-side sub-nodes N7b-0/1/2/3 and the glue N7a remain valid lemmas ---
they are reused by the Case~I producer, \cref{lem:case-I-realization}, which has no
single re-attached body and genuinely needs the row-counting/device route.) Build
order: M1 $\to$ M2 $\to$ M3 $\to$ N7. The motion-side route warrants its own
recon-before-build pass on M2 (the pin); M1 and M3 are mechanical.

\begin{lemma}[N7b-4 (superseded, row-side): recover $e_0$'s constraint from $v$'s two new edges; KT eq.~(6.12)]
  \label{lem:case-II-placement-e0-recovery}
  \uses{def:panel-hinge-framework, lem:case-II,
        lem:case-II-placement-old-rows-extract, lem:case-II-placement-old-rows,
        lem:case-II-placement-new-rows}
  \emph{Superseded by the motion-side re-frame above (M1--M3); kept for the audit
  trail.} This row-side formulation was found geometrically unbuildable on
  2026-06-04: it asks for an $e_0$-free independent old block of $G_v^{ab}$ of full
  count $D(|V(G)|-2)$, but no such block exists --- the $e_0$-row is not redundant
  in $G_v^{ab}$, since $G - v$ (the $e_0$-free subgraph) is not rigid (see ``Why the
  row-side framing fails'' above). The genuine rank-lift content --- recovering
  $e_0$'s span-membership constraint from $v$'s two new edges along the path
  $a$--$v$--$b$ --- is correct, but it lives \emph{in $G$} (motion-side), not as a
  row block of $G_v^{ab}$ alone, and is carried by \cref{lem:case-II-placement-pin-vertex}
  (M2). \emph{Red, superseded} (Phase~21b realization layer).
\end{lemma}
\begin{proof}
  Superseded; see the motion-side rank-lift (M1--M3) above.
\end{proof}

\begin{lemma}[M1: independent extensors have disjoint $1$-dimensional spans]
  \label{lem:case-II-placement-disjoint-line-meet}
  \uses{def:hinge-constraint}
  For $C, C' \in \mathrm{ScrewSpace}\,k$ linearly independent (in particular both
  nonzero), the $1$-dimensional spans meet trivially:
  $\operatorname{span}\{C\} \cap \operatorname{span}\{C'\} = \{0\}$. This is the meet
  step of the Case~II pin (\cref{lem:case-II-placement-pin-vertex}): a vector lying
  in both hinge spans is a scalar multiple of $C$ \emph{and} of $C'$, which forces it
  to vanish by independence. \emph{Red} (Phase~21b); pure linear algebra.
\end{lemma}
\begin{proof}
  A vector in the intersection is $\lambda C = \mu C'$; independence of $C, C'$
  forces $\lambda = \mu = 0$.
\end{proof}

\begin{lemma}[M2: two transversal hinges with independent extensors pin a re-attached body]
  \label{lem:case-II-placement-pin-vertex}
  \uses{def:hinge-constraint, def:panel-hinge-framework,
        lem:case-II-placement-disjoint-line-meet,
        lem:exists-independent-panel-extensor}
  The crux of the $1$-extension rank-lift (KT eq.~(6.12)), motion-side. Let $F$ be a
  body-hinge framework in which the body $v$ is joined to $a, b$ by hinges $e_a, e_b$
  whose support extensors $C(e_a), C(e_b)$ are linearly independent (the
  general-position condition, $m = 2$ case of
  \cref{lem:exists-independent-panel-extensor}). Then any infinitesimal motion $S$ of
  $F$ with $S a = S b$ also has $S v = S a$. Hence a motion constant on a body set
  containing $a, b$ is constant on that set together with $v$. Proof: the hinge
  constraints give $S v - S a \in \operatorname{span} C(e_a)$ and
  $S v - S b = S v - S a \in \operatorname{span} C(e_b)$, so the difference lies in
  the trivial meet (\cref{lem:case-II-placement-disjoint-line-meet}) and vanishes.
  \emph{Red} (Phase~21b); recon-before-build.
\end{lemma}
\begin{proof}
  \uses{lem:case-II-placement-disjoint-line-meet}
  From $S a = S b$ and the two hinge constraints,
  $S v - S a \in \operatorname{span} C(e_a) \cap \operatorname{span} C(e_b) = \{0\}$
  by \cref{lem:case-II-placement-disjoint-line-meet}, so $S v = S a$.
\end{proof}

\begin{lemma}[M3: Case II realization, motion-side]
  \label{lem:case-II-placement-motion-side-assembly}
  \uses{def:panel-hinge-framework, def:rank-hypothesis, lem:reduction-step,
        lem:case-II-placement-pin-vertex, lem:motions-mono-of-graph-le,
        lem:exists-independent-panel-extensor}
  The motion-side assembly producing $\mathrm{HasFullRankRealization}\,k\,G$ from the
  inductive realization of $G_v^{ab}$, superseding the row-side N7b path for Case~II.
  Choose $v$'s panel normal so $C(e_a), C(e_b)$ are linearly independent
  (\cref{lem:exists-independent-panel-extensor}, $m = 2$), placing the rest of $G$ at
  the inductive normals. The subgraph $G'' = (G-v)\cup\{e_a, e_b\} \le G$ shares
  $G - v$ with $G_v^{ab}$; on $G''$, a motion constant on $V(G)\setminus\{v\}$ (the
  inductive rigidity of $G_v^{ab}$, transported through $G - v$) is pinned at $v$ by
  \cref{lem:case-II-placement-pin-vertex}, so $G''$ is rigid on $V(G)$; re-adding $G$'s
  remaining edges keeps it rigid (\cref{lem:motions-mono-of-graph-le}). Feeds the
  resulting realization to \cref{lem:case-II-realization}. \emph{Red} (Phase~21b);
  the heart is M2.
\end{lemma}
\begin{proof}
  \uses{lem:case-II-placement-pin-vertex, lem:motions-mono-of-graph-le}
  Assemble M2 with the subgraph monotonicity
  \cref{lem:motions-mono-of-graph-le} and the inductive rigidity on
  $V(G)\setminus\{v\}$; see \cite[\S6.3]{katohTanigawa2011} (Lemma~6.8, eq.~(6.12)).
\end{proof}


\subsection{Case III deferred (the crux)}
\label{sec:molecular-algebraic-induction-caseIII}

The crux of Case III is the missing $+1$ row of KT Claim~6.11 (the $D$-th
$ab$-row of $R(G_v^{ab}, q)$ that lifts the stratum-1 brick
$D(|V|-1) - 1$ of \cref{lem:case-II-realization-placement} to full rank
$D(|V|-1)$). Claim~6.11's discharge factors, in dependency order, through three
combinatorial$\to$linear gaps; the leaf-most --- pure $M(\tilde G)$ matroid
theory, all inputs green --- is the matroid-base form of KT Lemma~4.3(ii) at
$k = 0$, built here.

\begin{lemma}[Matroid-base 4.3(ii) at $k = 0$; KT Lemma~4.3(ii), Claim~6.11 prerequisite]
  \label{lem:case-III-claim-6-11-base}
  \lean{Graph.splitOff_exists_base_inter_fiber_lt}
  \leanok
  \uses{def:k-dof, def:matroid-MG, thm:def-eq-corank, lem:forest-surgery-count, lem:splitoff-deficiency}
  Let $G$ be a $0$-dof-graph ($\operatorname{def}(\tilde G) = 0$) with a degree-2
  vertex $v$ (distinct neighbours $a, b$, incident edges exactly $e_a \ne e_b$,
  $e_0 \notin E(G)$ fresh). Then $M(\widetilde{G_v^{ab}})$ has a base $B'$ with
  \[
    \bigl|\widetilde{ab} \cap B'\bigr| \;<\; D - 1,
  \]
  i.e.\ $\widetilde{ab} \not\subseteq B'$: a redundant $\widetilde{ab}$-copy exists.
  This is KT's step~1 toward Claim~6.11 (the redundant $ab$-row of
  $R(G_v^{ab}, q)$); the remaining gaps --- the nested IH at the restricted
  realization and the $M(\tilde G)\leftrightarrow$row-dependence bridge --- stay
  deferred.
\end{lemma}
\begin{proof}
  \leanok
  Reroute a balanced $D$-forest packing of a base of $M(\tilde G)$ across $v$
  (\cref{lem:forest-surgery-count}): this yields an $M(\widetilde{G_v^{ab}})$-%
  independent set $I'$ with $|I'| = |\text{base}| - D$ and $|\widetilde{ab} \cap I'|
  < D - 1$ (KT Lemma~4.1's two conclusions). At $k = 0$ splitting off keeps the
  graph $0$-dof (\cref{lem:splitoff-deficiency}, with
  $\operatorname{def} \ge 0$), so $\operatorname{rank} M(\widetilde{G_v^{ab}}) =
  D(|V \setminus \{v\}| - 1) = |\text{base}| - D = |I'|$ by the def\,$=$\,corank
  identity (\cref{thm:def-eq-corank}); an independent set of full rank is a base.
  That base $I'$ carries the fiber bound.
\end{proof}

The second factor of Claim~6.11 is KT's nested-induction step: the
vertex-removal $G_v := G - v = G_v^{ab} - ab$ is itself a \emph{minimal}
$k'$-dof-graph for $k' := \operatorname{def}(\tilde G_v)$, and that deficiency is
small ($k' \le D - 2$). The combinatorial half --- minimality and the deficiency
bound, both pure $M(\tilde G)$ matroid theory --- is green here; the
\emph{analytic} half it feeds (applying the geometric IH to $G_v$ at the
inductively-fixed seed, KT eq.\ (6.22)) is the deferred crux.

\begin{lemma}[$G - v$ is minimal $k'$-dof with $k' \le D - 2$; KT Lemma~4.3(ii)/nested IH, Claim~6.11 prerequisite]
  \label{lem:case-III-gap3-minimalKDof}
  \lean{Graph.splitOff_removeVertex_minimalKDof}
  \leanok
  \uses{def:k-dof, def:graph-operations, lem:subgraph-minimality, thm:def-eq-corank,
        lem:splitoff-deficiency, lem:case-III-claim-6-11-base}
  Let $G$ be a minimal $0$-dof-graph with a degree-2 vertex $v$ (distinct
  neighbours $a, b$, incident edges exactly $e_a \ne e_b$, $e_0 \notin E(G)$
  fresh). Then the vertex-removal $G_v := G - v$ is a minimal
  $k'$-dof-graph for $k' = \operatorname{def}(\tilde G_v)$, with
  \[
    0 \;\le\; k' \;\le\; D - 2.
  \]
\end{lemma}
\begin{proof}
  \leanok
  Minimality transports along $G_v \le G$ by \cref{lem:subgraph-minimality}: a
  subgraph of a minimal $k$-dof-graph is minimal at its own deficiency. For the
  bound, take the Gap-2 base $B'$ of $M(\widetilde{G_v^{ab}})$ with
  $h := |\widetilde{ab} \cap B'| < D - 1$
  (\cref{lem:case-III-claim-6-11-base}). Its surviving part $B' \setminus
  \widetilde{ab}$ lands in $E(\tilde G_v)$ and is independent in $M(\tilde G_v) =
  M(\widetilde{G_v^{ab}}) \restriction E(\tilde G_v)$ (matroid restriction along
  $\tilde G_v \le \widetilde{G_v^{ab}}$), so $\operatorname{rank} M(\tilde G_v)
  \ge |B'| - h$. At $k = 0$ the splitting off is $0$-dof
  (\cref{lem:splitoff-deficiency}), so $|B'| = D(|V \setminus \{v\}| - 1)$; the
  def\,$=$\,corank identity (\cref{thm:def-eq-corank}, same vertex set $V(G)
  \setminus \{v\}$) then gives $\operatorname{def}(\tilde G_v) = D(|V \setminus
  \{v\}| - 1) - \operatorname{rank} M(\tilde G_v) \le h < D - 1$, i.e.\
  $k' \le D - 2$. The lower bound $0 \le k'$ is the nonnegativity of deficiency.
\end{proof}

\begin{lemma}[Case III: $k = 0$, no proper rigid subgraph; KT Lemma 6.10/6.13]
  \label{lem:case-III}
  \uses{thm:minimal-kdof-reduction, lem:extensor-independence,
        lem:cycle-realization, lem:case-II-realization-placement,
        lem:case-III-claim-6-11-base, lem:case-III-gap3-minimalKDof}
  Let $G$ be a minimal $0$-dof-graph that is $2$-edge-connected with no
  proper rigid subgraph. Then $G$ has a panel-hinge realization at rank
  $D(|V|-1)$. A single candidate realization is one row short of full
  rank (the stratum-1 brick $D(|V|-1) - 1$ of
  \cref{lem:case-II-realization-placement}); the missing $+1$ row is KT
  Claim~6.11, whose combinatorial prerequisites --- the matroid-base 4.3(ii)
  form (\cref{lem:case-III-claim-6-11-base}) and the nested-IH reduction
  $G - v$ minimal $k'$-dof with $k' \le D - 2$
  (\cref{lem:case-III-gap3-minimalKDof}) --- are green; what stays deferred
  is the analytic eq.\ (6.22) rank-at-the-fixed-seed transfer. KT then
  build $D$ candidate realizations and show one is rigid; the
  degree-2 condition forces all candidates to test the same normal vector $r
  \in \mathbb{R}^D$, which by the extensor-independence Lemma~2.1
  (\cref{lem:extensor-independence}) must vanish --- contradiction.
  \textbf{Deferred to Phases~22--23} ($d = 3$ first, then general $d$); stated
  here only to close the Theorem~5.5 dichotomy.
\end{lemma}
\begin{proof}
  Deferred to Phases~22--23. See \cite[\S6.4]{katohTanigawa2011}.
\end{proof}

